     1	\ifx\danslelivre\undefined
     2	\RequirePackage{luatex85}
     3	\documentclass[10pt]{../fichiers-configuration-latex/smfart}
     4	\input{../fichiers-configuration-latex/paquets}
     5	\input{../fichiers-configuration-latex/numerotation}
     6	\input{../fichiers-configuration-latex/macros}
     7	
     8	
     9	\externaldocument{../01/01.tex}
    10	\externaldocument{../02/02.tex}
    11	\externaldocument{../03/03.tex}
    12	\externaldocument{../04/04.tex}
    13	\externaldocument{../05/05.tex}
    14	\externaldocument{../05bis/05bis.tex}
    15	\externaldocument{../06/06.tex}
    16	\externaldocument{../07/07.tex}
    17	\externaldocument{../08/08.tex}
    18	%\externaldocument{../09/09.tex}
    19	\externaldocument{../10/10.tex}
    20	\externaldocument{../11/11.tex}
    21	\externaldocument{../12/12.tex}
    22	\externaldocument{../13/13.tex}
    23	\externaldocument{../14/14.tex}
    24	\externaldocument{../15/15.tex}
    25	\externaldocument{../16/16.tex}
    26	\externaldocument{../17/17.tex}
    27	\externaldocument{../18/18.tex}
    28	\externaldocument{../19/19.tex}
    29	
    30	\setboolean{original}{false}
    31	
    32	\begin{document}
    33	
    34	\begin{center}
    35	IX. Faisceaux constructibles cohomologie d'une courbe algébrique\\
    36	M. Artin 
    37	\end{center}
    38	
    39	\selectlanguage{english}version : \showgitstatus\selectlanguage{french}
    40	\tableofcontents
    41	
    42	\else
    43	\chapter{Faisceaux constructibles cohomologie d'une courbe algébrique}
    44	%\addtocontents{toc}{M. Artin}
    45	\begin{center}
    46	M. Artin 
    47	\end{center}
    48	\fi
    49	
    50	\setcounter{pageoriginal}{0}
    51	
    52	\setcounter{section}{-1}
    53	\section{Introduction}\etiquette[0]{IX.sec0}\pageoriginale
    54	
    55	On notera qu'à partir du présent exposé, contrairement aux exposés
    56	précédents et à la théorie cohomologique classique des espaces
    57	topologiques, nous sommes obligés, pour la validité des énoncés
    58	essentiels, de nous limiter aux faisceaux \emph{de torsion} (et
    59	souvent même, plus précisément, aux faisceaux de torsion premiers aux
    60	caractéristiques résiduelles).
    61	
    62	Signalons qu'en fait une théorie « raisonnable »  de la
    63	cohomologie à coefficients entiers (ou même réels), pour les variétés
    64	sur un corps algébriquement clos de caractéristique $p\neq 0$,
    65	analogue à la théorie classique pour le cas
    66	$k=\CC$, n'existe
    67	pas (comme l'a remarqué Serre). Plus précisément, il n'existe pas de
    68	foncteur contravariant $H^1$, défini disons sur la catégorie des
    69	schémas projectifs lisses sur le corps $K$ alg. clos de car. $p> 0$, à
    70	valeurs dans la catégorie des espaces vectoriels de dimension finie
    71	sur $\RR$  (ou un sous-corps de $\RR$),
    72	« commutant aux produits », \ie tel que $H^1(X\times
    73	Y)\xleftarrow{\sim}H^1(X)\times H^1(Y)$, et tel que $\dim H^1(X)=2$ si
    74	$X$ est une courbe connexe de genre 1. En effet, il s'ensuivrait que
    75	si $X$ est une variété abélienne sur $k$, alors l'application
    76	$u\mapsto u^1$ de $\End(X)$ dans $\End(H^1(X))$ est \emph{additive},
    77	donc une \emph{représentation} de l'anneau opposé de $\End(X)$. Mais
    78	en caractéristique $p$, il existe des courbes
    79	elliptiques\pageoriginale $X$ ayant comme anneau d'endomorphismes un
    80	ordre maximal $A$ d'une algèbre de quaternions définie sur
    81	$\ZZ$ ([3], p.198), et une telle algèbre n'a
    82	évidemment pas de représentation dans un espace vectoriel de dimension
    83	2 sur $\RR$, puisque
    84	$A\otimes_{\ZZ}\RR$ est un corps
    85	(le corps des quaternions).
    86	
    87	\section{Le sorite des faisceaux de torsion}\etiquette[1]{IX.sec1}
    88	
    89	Soient $T$ un topos et $F$ un faisceau abélien sur $T$. On peut
    90	multiplier par $n$ dans $F$ ($n\in \ZZ$ un nombre
    91	entier), cette multiplication étant induite par l'opération analogue
    92	dans (Ab). Notons ${}_nF$ le noyau de cette multiplication; donc
    93	${}_nF(X)$, pour $X\in \Ob T$, est le groupe des sections de $F(X)$
    94	dont l'ordre divise $n$.
    95	
    96	On désigne par $\PP$ l'ensemble de tous les nombres
    97	premiers.
    98	
    99	\begin{enonce*}{Définition 1.1}\etiquette[1.1]{IX.def1.1}
   100	Soit $p$ un ensemble de nombres premiers. On dit que $F$ est un
   101	faisceau de $p$-torsion, \sisi{où}{ou} un $p$-faisceau, si le morphisme canonique
   102	$$
   103	\varinjlim_n {}_nF\longrightarrow F\text{ est bijectif}{\quoi},
   104	$$
   105	$n$ parcourant l'ensemble des entiers tels que $\ass n\subset p$,
   106	c'est-à-dire, tels que les nombres premiers divisant $n$ soient dans
   107	$p$. Si $p=\PP$ est l'ensemble de tous les nombres
   108	premiers, on dit simplement que $F$ est de torsion.
   109	\end{enonce*}
   110	
   111	\begin{enonce*}{Proposition 1.2}\etiquette[1.2]{IX.prop1.2}
   112	\begin{enumerate}
   113	\renewcommand{\theenumi}{\roman{enumi}}
   114	\renewcommand{\labelenumi}{(\theenumi)}
   115	\item $F$ est de $p$-torsion si et seulement si $F$ est le faisceau
   116	  associé à un préfaisceau à valeurs dans des groupes abéliens de
   117	  $p$-torsion; on peut prendre
   118	$$
   119	P=\frac{\lim(\text{\rm préf})}{\ass n\subset p} {}_nF{\quoi}.
   120	$$
   121	
   122	\item Si\pageoriginale $F$ est de $p$-torsion et si $X\in \Ob T$ est
   123	  un objet quasi-compact\footnote[1]{Un objet $X$ d'un topos est dit
   124	  quasi-compact si chaque famille couvrante $\{X_i\rightarrow X\}$
   125	  peut être majorée par une famille couvrante finie. Un morphisme $u$
   126	  de topos est quasi-compact si l'image inverse de chaque objet
   127	  quasi-compact est encore quasi-compacte.}, alors $F(X)$ est un
   128	  groupe de $p$-torsion.
   129	
   130	\item Si $T$ est localement de type fini {\rm (VI 1.1)}, alors $F$ est
   131	  de $p$-torsion si et seulement si $F(X)$ est un groupe de
   132	  $p$-torsion pour chaque $X\in \Ob T$ quasi-compact. Dans ce cas, on
   133	  a
   134	$$
   135	\varinjlim_{\ass n\subset p} H^q(T/X; {}_nF)\xrightarrow{\sim}
   136	H^q(T/X;F)
   137	$$
   138	pour chaque $X$ quasi-compact et tout $q$.
   139	
   140	\item Si $u:T\rightarrow T'$ est un morphisme de topos et si $F'$ est
   141	  de $p$-torsion sur $T'$, alors l'image inverse $u^{\ast}F$ est un
   142	  faisceau de $p$-torsion sur $T$.
   143	
   144	\item Si $u:T\rightarrow T'$ est un morphisme de topos, avec $T$ et
   145	  $T'$ localement de type fini et $u$ quasi-compact${}^\ast$, et si $F$ est un faisceau de $p$-torsion sur $T$, alors
   146	  les $R^qu_{\ast}F'$ sont des faisceaux de $p$-torsion sur $T$ pour
   147	  tout $q$.
   148	\end{enumerate}
   149	\end{enonce*}
   150	
   151	\noindent
   152	\emph{Démonstration}:
   153	\begin{enumerate}
   154	\renewcommand{\theenumi}{\roman{enumi}}
   155	\renewcommand{\labelenumi}{(\theenumi)}
   156	\item Évidemment, si un faisceau $F$ est de $p$-torsion, $F$ est le
   157	  faisceau associé au préfaisceau $P$ dont le groupe des sections
   158	  $P(X)$ est le sous-groupe des éléments de $p$-torsion de
   159	  $F(X)$. Inversement, soit $F=\underline{a}P$ (II) le faisceau associé à $P$,
   160	  où $P$ est à valeurs dans la catégorie des groupes abéliens de
   161	  $p$-torsion. De la suite exacte
   162	$$
   163	0\rightarrow {}_nP \rightarrow P\xrightarrow{n}P
   164	$$
   165	on déduit une suite exacte
   166	$$
   167	0\rightarrow \underline{a}({}_nP)\rightarrow F\xrightarrow{n}F{\quoi},
   168	$$
   169	d'où\pageoriginale $\underline{a}({}_nP)={}_nF$. Comme le foncteur $\underline{a}$
   170	commute aux limites inductives, on déduit de
   171	$$
   172	\frac{\lim(\text{\rm préf})}{\ass n\subset p}{}_nP\xrightarrow{\sim}P
   173	$$
   174	qu'on a aussi
   175	$$
   176	\varinjlim_{\ass n\subset p} {}_nF\xrightarrow{\sim}F{\quoi}.
   177	$$
   178	
   179	\item Posons $P=\underrightarrow{\lim(\text{\rm préf})}{}_n F$. Alors on a
   180	  $P\subset F$, donc $P$ est un préfaisceau séparé, et par suite
   181	$$
   182	\underline{a}P(X)=\varinjlim (\ker(\prod_i P(X_i)\rightrightarrows)){\quoi},
   183	$$
   184	où comme d'habitude la limite est prise suivant les familles
   185	couvrantes $\{X_i\rightarrow X\}$. Mais comme $X$ est quasi-compact,
   186	il suffit de prendre les familles couvrantes finies. Pour ces familles
   187	$\prod_iP(X_i)$ est de $p$-torsion; donc le $\ker$ l'est aussi. Par
   188	suite $\underline{a}P(X)=F(X)$ est de $p$-torsion.
   189	
   190	\item Pour vérifier que $\varinjlim {}_nF\rightarrow F$ est bijectif,
   191	  il suffit de le faire pour les valeurs sur $X\in \Ob T$ pour chaque
   192	  $X$ d'une famille de générateurs. On est donc ramené à (ii) pour la
   193	  première assertion. La deuxième assertion est une conséquence de
   194	
   195	\item Conséquence de (i) puisque $u^{\ast}(\underline{a}P')=\underline{a}(u\cdot P)$
   196	
   197	\item Conséquence facile de (iii).
   198	\end{enumerate}
   199	
   200	On aura aussi besoin d'une variante non-abélienne. Mais il y a des
   201	problèmes techniques dans la définition (j'ignore s'ils sont sérieux),
   202	et on se contentera donc de donner une définition pour la topologie
   203	étale d'un schéma.
   204	
   205	\begin{enonce*}{Définition 1.3}\etiquette[1.3]{IX.def1.3}
   206	Soit\pageoriginale $p$ un ensemble de nombres premiers. On dit qu'un
   207	groupe $G$ est un \emph{ind-p-groupe} si chaque sous-ensemble fini de $G$
   208	engendre un sous-groupe fini d'ordre $n$, avec $\ass n\subset p$. Il
   209	revient au même de dire que les sous-groupes finie $G_{\alpha}$
   210	d'ordre $n_{\alpha}$, avec $\ass n_{\alpha}\subset p$, forment un
   211	ensemble filtrant, et que $G=\varinjlim G_{\alpha}$. Si
   212	$p=\PP$, on dit \emph{groupe ind-fini} au lieu de ind-p-groupe.
   213	\end{enonce*}
   214	
   215	On vérifie immédiatement la
   216	\begin{enonce*}{Proposition 1.4}\etiquette[1.4]{IX.prop1.4}
   217	\begin{enumerate}
   218	\renewcommand{\theenumi}{\roman{enumi}}
   219	\renewcommand{\labelenumi}{(\theenumi)}
   220	\item Un sous-groupe ainsi qu'un quotient d'un ind $p$-groupe est
   221	  encore un ind $p$-groupe.
   222	
   223	\item Une limite inductive filtrante de ind-$p$-groupes est un
   224	  ind-$p$-groupe.
   225	
   226	\item Une limite projective finie de ind-p-groupes est un ind-p-groupe.
   227	\end{enumerate}
   228	\end{enonce*}
   229	
   230	\begin{enonce*}{Définition 1.5}\etiquette[1.5]{IX.def1.5}
   231	Soit $X$ un schéma. On dit qu'un faisceau $F$ de groupes sur $X$ est
   232	un \emph{faisceau de ind $p$-groupes} (\resp de \emph{groupes
   233	  ind-finis})
   234	si pour chaque $U\rightarrow X$ étale, avec $U$ quasi-compact, $F(U)$ est un
   235	ind $p$-groupe (\resp un groupe ind-fini).
   236	\end{enonce*}
   237	
   238	\begin{enonce*}{Proposition 1.6}\etiquette[1.6]{IX.prop1.6}
   239	Soit $F$ un faisceau de groupes sur le schéma $X$.
   240	\begin{enumerate}
   241	\renewcommand{\theenumi}{\roman{enumi}}
   242	\renewcommand{\labelenumi}{(\theenumi)}
   243	\item $F$ est un faisceau de ind-$p$-groupes si et seulement si pour
   244	  chaque point géométrique $\xi$ de $X$, la fibre $F_{\xi}$ est un
   245	  ind-p-groupe.
   246	
   247	\item Si $f:X\rightarrow X'$ est un morphisme et si $F'$ est un
   248	  faisceau de ind-p-groupes sur $X'$, $f^{\ast}F'$ est un faisceau de
   249	  ind-p-groupes.
   250	
   251	\item Si $f:X\rightarrow X'$ est un morphisme quasi-compact et $F$ est
   252	  un faisceau de ind-p-groupes sur $X$, alors $f_{\ast}F$ est un
   253	  faisceau de ind-p-groupes.
   254	\end{enumerate}
   255	\end{enonce*}
   256	
   257	\noindent
   258	\emph{Démonstration.}
   259	(i) Supposons que $F$ soit un faisceau de ind-p-groupes, et soit
   260	  $\xi$ un point géométrique de $X$. Alors $F=\varinjlim F(X')$, où
   261	  $X'$ parcourt un\pageoriginale système pseudo-filtrant de schémas
   262	étales sur $X$ (VIII 3.9). Il est évident que les $X'$ affines (donc
   263	quasi-compacts) forment un ensemble cofinal, et que par suite $F$ est
   264	limite inductive de ind $p$-groupes, donc est un ind-p-groupe d'après
   265	\hyperref[IX.prop1.4]{1.4} (ii). Inversement, supposons que pour chaque $\xi$ la fibre
   266	$F_{\xi}$ soit un ind-p-groupe et soit $U\rightarrow X$ étale, $U$
   267	quasi-compact. Soit $S\subset F(U)$ un sous-ensemble fini. Pour chaque
   268	point géométrique $\xi$ de $U$, l'image de $S$ dans $F_{\xi}$ engendre
   269	un groupe fini. Comme un groupe fini est de présentation finie, on
   270	conclut aisément qu'il existe un $U_{\xi}$ étale sur $U$,
   271	$\xi$-ponctué, tel que $S$ engendre un groupe fini dans $F(U_{\xi})$
   272	(\cf VIII 4). Un nombre fini de tels $U_{\xi}$ forment un recouvrement
   273	de $U$, soit $\{U_1,\ldots,U_n\}$. Alors l'image de $S$ dans $\prod
   274	F(U_i)$ engendre un sous-groupe fini, et comme $F(U)\subset \prod
   275	F(U_i)$, on a gagné.
   276	
   277	L'assertion (ii) est triviale à partir de (i), et (iii) est immédiate.
   278	
   279	\section{Faisceaux constructibles}\etiquette[2]{IX.sec2}
   280	
   281	\setcounter{subsection}{-1}
   282	\subsection{}\etiquette[2.0]{IX.subsec2.0}
   283	Soit $T$ un topos. Rappelons qu'à chaque élément $S\in \Ob (\Ens)$ on
   284	associe le faisceau $S_T$ associé au préfaisceau $P$ tel que $P(X)=S$
   285	pour chaque $X\in \Ob (T)$. L'objet de $T$ qui représente $S_T$ est
   286	$$
   287	\coprod_{s\in S}e= ``S\times e \mbox{''}{\quoi},
   288	$$
   289	$e$ l'objet final de $T$. $S_T$ est appelé \emph{faisceau constant},
   290	de valeur $S$ (IV). On définit d'une manière évidente la notion de
   291	faisceau \emph{de groupes constant, de $A$-modules constant} ($A$ un
   292	anneau), et de faisceau \emph{abélien constant}, de valeur $M$. Un
   293	morphisme $f_T:S_T\rightarrow S'_T$ de faisceaux constants est dit
   294	constant s'il provient d'un morphisme d'ensembles $f:S\rightarrow S'$.
   295	
   296	Un faisceau $F$ est dit \emph{localement constant} s'il existe un
   297	recouvrement $\{e_i\rightarrow e\}$ de l'objet final tel que $F$
   298	devienne constant sur chaque $e_i$. On\pageoriginale définit d'une
   299	façon analogue la notion de faisceau \emph{de groupes} ou de
   300	$A$-\emph{modules localement constant}, et de \emph{morphisme
   301	  localement constant} $f:F\rightarrow G$, si $F$ et $G$ sont des
   302	faisceaux localement constants. Enfin, un faisceau de groupes ou de
   303	$A$-modules localement constant est dit \emph{de type fini} (\resp
   304	\emph{de présentation finie}) si les valeurs locales sont des groupes
   305	ou des $A$-modules de type fini (\resp de présentation finie).
   306	
   307	\begin{enonce*}{Lemme 2.1}\etiquette[2.1]{IX.lem2.1}
   308	\begin{enumerate}
   309	\renewcommand{\theenumi}{\roman{enumi}}
   310	\renewcommand{\labelenumi}{(\theenumi)}
   311	\item Soit $f:F\rightarrow G$ un morphisme de faisceaux d'ensembles
   312	  localement constant sur $T$, et supposons que les valeurs locales de
   313	  $F$ soient des ensembles finis. Alors $f$ est localement constant.
   314	
   315	\item Soit $f:F\rightarrow G$ un morphisme de faisceaux de $A$-modules
   316	  localement constants, avec $F$ de type fini. Alors $f$ est
   317	  localement constant, et le noyau et conoyau de $f$ sont localement
   318	  constants.
   319	
   320	\item Soit $0\rightarrow F'\rightarrow F\rightarrow F''\rightarrow 0$
   321	  une suite exacte de faisceaux de $A$-modules (\resp de groupes),
   322	  avec $F'$ et $F''$ localement constants, $F''$ étant de présentation
   323	  finie. Alors $F$ est localement constant.
   324	\end{enumerate}
   325	\end{enonce*}
   326	
   327	Démonstration de (i) et (ii). On se ramène sans peine au cas où $F$ et
   328	$G$ sont constants. Traitons (ii): soit $F=M_T$, $G=N_T$ ($M$, $N$ des
   329	$A$-modules), et soit $S\subset M$ un sous-ensemble fini de
   330	générateurs de $M$. Alors $f$ est déterminé par sa restriction à $S_T$
   331	et $f$ sera localement constant si $f|S_T$ l'est. On est donc
   332	réduit à (i) pour la première assertion de (ii), et la deuxième
   333	assertion est conséquence triviale de la première.
   334	
   335	Il reste donc à démontrer (i). Supposons que $F=S_T$, $G={S'}_T$
   336	soient constants, et notons aussi par $f$ le morphisme d'objets
   337	$f:S\times e\rightarrow S'\times e$ (notation comme ci-dessus). Soient
   338	$f_s:e\rightarrow S'\times e$ $(s'\in S)$ les composants de $f$: puisque
   339	$S$ est fini, il suffit évidemment de traiter le cas $f=f_s$,
   340	c'est-à-dire, où $S$ est un ensemble d'un élément. Soit
   341	$X_{s'}\rightarrow e$ $(s'\in S')$ la famille des morphismes qui rendent
   342	cartésiens les diagrammes
   343	$$
   344	\xymatrix@C=2cm@R=1.5cm{
   345	X_{s'} \ar[r] \ar[d]& e\ar[d]^{i_{S'}}\\
   346	e   \ar[r]^{f}   & S'\times e
   347	}
   348	$$
   349	où\pageoriginale $i_{s'}:e\rightarrow S'\times e$ est
   350	« l'inclusion » dans la $s'$--ième composante. La famille $\{i_{s'}\}$ est
   351	trivialement couvrante, donc $\{X_{s'}\rightarrow e\}$ l'est aussi. On
   352	vérifie immédiatement (puisque les sommes directes dans un topos sont
   353	disjointes) %% référence manquante
   354	que $f$ devient constant sur $X_{s'}$, et que par suite $f$
   355	est localement constant.
   356	
   357	Démonstration de (iii). Traitons le cas des faisceaux de
   358	$A$-modules. On peut supposer $F'$ et $F''$ constants, $F''$ de valeur
   359	$M''$ un $A$-module de présentation finie. Si $M''$ est libre, donc
   360	admet une base finie $(e_i)_{1\leq i\leq n}$, alors les $e_i$
   361	définissent des sections de $F''$, donc se relèvent localement en des
   362	sections de $F$, ce qui montre que localement l'extension $F$ de $F''$
   363	par $F'$ se scinde, ce qui prouve que $F$ est localement isomorphe à
   364	$F'\times F''$, donc est localement constant. Dans le cas général,
   365	choisissant un homomorphisme surjectif $\varphi:L''\rightarrow M''$,
   366	avec $L''$ libre de type fini; le noyau $R$ de $\varphi$ est de type
   367	fini, et posant $L=\fprod{L''_T}{F}{M''_T}=\fprod{L''_T}{F}{F''}$; on
   368	voit d'une part que $L$ est une extension de $L''_T$ par $G$, donc
   369	localement constant d'après ce qui précède, d'autre part qu'on a une
   370	suite exacte $0\rightarrow R_T\rightarrow L\rightarrow F\rightarrow
   371	0$, ce qui grâce à (ii) implique que $F$ est localement constant.
   372	
   373	A partir de maintenant, on se borne aux topos étales des schémas.
   374	
   375	\begin{enonce*}{Lemme 2.2}\etiquette[2.2]{IX.lem2.2}
   376	Soit $X$ un schéma et $F$ un faisceau sur $X$ qui est localement
   377	constant. Alors $F$ est représenté par un $Y/X$ étale. Si de plus les
   378	fibres de $F$ sont finies (\resp et non-vides), alors $Y$ est un
   379	revêtement étale (\resp et surjectif).
   380	\end{enonce*}
   381	
   382	Démonstration. Descente (\SGA 1 IX 4.1 si $F$ à fibres finies, \SGA 3 X
   383	5.4 dans le cas général).
   384	
   385	\begin{enonce*}{Définition 2.3}\etiquette[2.3]{IX.def2.3}
   386	Un faisceau d'ensembles (\resp de groupes, \resp de $A$-modules) est
   387	dit constructible si pour chaque ouvert affine $U\subset X$ il existe
   388	une décomposition de $U$ en réunion d'un nombre fini de sous-schémas
   389	constructibles {\rm (\EGA $0_{{\rm III}}$ 9.1.2)} localement fermés
   390	réduits $U_i$ telle que le faisceau induit par $F$\pageoriginale sur
   391	chaque $U_i$ soit localement constant, de valeur finie (\resp finie,
   392	\resp de présentation finie\footnote[1]{Lorsque $A$ n'est pas
   393	  supposé noethérien, il est préférable, pour les besoins de l'Algèbre
   394	Homologique, d'exiger ici, au lieu de la seule présentation finie pour
   395	le $A$-module $M$, qu'il ait une résolution gauche par des $A$-modules
   396	libres de type fini. La notion introduite dans (\hyperref[IX.def2.3]{2.3}) devrait alors
   397	prendre le nom: $F$ est 1-constructible, (plus généralement, on
   398	définirait de façon évidente la $n$-constructibilité, pour tout
   399	entier $n\geq 0$). Nous allons provisoirement pour le présent exposé
   400	garder la terminologie sous la forme (\hyperref[IX.def2.3]{2.3}), qui est surtout
   401	raisonnable dans le cas où $A$ est noethérien, où elle coïncide avec
   402	celle qu'on vient de signaler. Comparer \SGA 6 I pour des notions de
   403	finitude dans des cas non noethériens.}.
   404	\end{enonce*}
   405	
   406	\setcounter{subsection}{3}
   407	\subsubsection{}\etiquette[2.3.1]{IX.subsubsec2.3.1}
   408	L'hypothèse de finitude sur les valeurs locales revient à dire que
   409	pour chaque point géométrique $\xi$, la fibre $F_{\xi}$ est finie
   410	(\resp finie, \resp de présentation finie). Il s'ensuit immédiatement
   411	de \hyperref[IX.lem2.1]{2.1} (i) qu'un faisceau de groupes $F$ est constructible si et
   412	seulement si le faisceau d'ensembles sous-jacent à $F$ est
   413	constructible.
   414	
   415	On fera attention que si $F$ est un faisceau abélien, il est
   416	constructible en tant que faisceau de groupes (ou encore, en tant que
   417	faisceau d'ensembles) si et seulement si il est constructible comme
   418	faisceau de $\ZZ$-modules, et si \emph{de plus} ses
   419	fibres sont finies. Ainsi, $\ZZ_X$ est
   420	constructible comme $\ZZ$-module, mais non comme
   421	faisceau de groupes.
   422	
   423	\begin{enonce*}{Proposition 2.4}\etiquette[2.4]{IX.prop2.4}
   424	\begin{enumerate}
   425	\renewcommand{\theenumi}{\roman{enumi}}
   426	\renewcommand{\labelenumi}{(\theenumi)}
   427	\item Soit $X$ quasi-compact et quasi-séparé. Alors un faisceau \\ (\resp
   428	  faisceau de groupes, \resp faisceau de $A$-modules) est
   429	  constructible si et seulement s'il existe une décomposition de $X$
   430	  en réunion de parties localement fermées constructibles, telle que
   431	  $F$ devienne localement constant, et fini (\resp fini, \resp de
   432	  présentation finie), sur chaque $X_i$.
   433	
   434	\item Pour vérifier que $F$ est constructible, il suffit de vérifier
   435	  que $F| U_i$ est constructible pour les $U_i$ d'un recouvrement
   436	  ouvert donné de $X$.
   437	
   438	\item Soit $F:Y\rightarrow X$ un morphisme. Alors $f^{\ast}F$ est
   439	  constructible si $F$ l'est.
   440	
   441	\item L'ensemble\pageoriginale des points où la fibre d'un faisceau
   442	  constructible est non vide (\resp non-nulle) est un sous-ensemble
   443	  localement constructible {\rm (\EGA $0_{\rm III}$ 9.1.2)}.
   444	
   445	\item Soit $X$ localement noethérien. Alors un faisceau d'ensembles
   446	  (\resp\ldots) $F$ sur $X$ est constructible si et seulement si pour
   447	  chaque $x\in X$, il existe un ouvert non-vide de l'adhérence
   448	  $\bar{x}$ de $x$ tel que $F$ induise un faisceau localement constant
   449	  de valeur finie (\resp\ldots) sur $U$.
   450	\end{enumerate}
   451	\end{enonce*}
   452	
   453	Démonstration. Évidemment, si $f:Y\rightarrow X$ est un morphisme et
   454	s'il existe une décomposition de $X$ en réunion de parties $X_i$
   455	localement fermées et constructibles, telles que $F$ devienne
   456	localement constant sur chaque $X_i$, il en est de même de $Y$ et
   457	$f^{\ast}F$ (\cf \EGA IV 1.8.2). Cela prouve l'implication $\Leftarrow$
   458	de (i). Supposons maintenant que $X$ soit quasi-compact et
   459	quasi-séparé et soit $\{U_i,\ldots,U_n\}$ un recouvrement ouvert
   460	affine de $X$ tel que $F|U_i$ soit constructible pour chaque
   461	$i$. Pour trouver une décomposition de $X$ comme dans (i), il suffit
   462	de le faire pour $U_n$ et pour $Y=X-U_n$ (avec la structure induite
   463	réduite), $U_n$ et $Y$ étant constructibles dans $X$ grâce à
   464	l'hypothèse « $X$ quasi-séparé ». Or une décomposition existe pour
   465	$U_n$ d'après la définition de la constructibilité, et $Y$ est réunion de
   466	$n-1$ ouverts affines $V_i=Y\cap U_i$ $(i=1,\ldots,n-1)$ avec $F|
   467	V_i$ constructible, d'où le résultat (i) par récurrence sur $n$. Le
   468	même raisonnement démontre (ii); en effet, l'hypothèse implique que
   469	$F| U$ est constructible pour chaque ouvert affine $U$ « assez
   470	petit », et on peut recouvrir un ouvert affine $V$ arbitraire par un
   471	nombre fini de tels ouverts. L'assertion (ii) montre que la notion de
   472	constructibilité est locale sur $X$, et (iii), (iv) s'ensuivent
   473	immédiatement. L'implication $\Rightarrow$ de (v) est immédiate, et ne
   474	dépend pas de l'hypothèse noethérienne. Si $X$ est noethérien on
   475	démontre l'implication $\Leftarrow$ par la « récurrence
   476	noethérienne » habituelle.
   477	
   478	\begin{enonce*}{Proposition 2.5}\etiquette[2.5]{IX.prop2.5}
   479	Soit $X$ quasi-compact et quasi-séparé, et $F$ un faisceau de groupes
   480	ou de $A$-modules constructible sur $X$. Alors il existe une
   481	filtration finie de $F$ dont les quotients successifs sont de la forme
   482	$i_{!}G$, où $i:U\rightarrow X$ est l'inclusion d'une partie
   483	localement fermée et constructible, et où $G$ est localement constant
   484	et constructible sur $U$. Si $X$ est noethérien, il existe une telle
   485	filtration, avec des $U$ irréductibles.
   486	\end{enonce*}
   487	
   488	Démonstration.\pageoriginale D'après \hyperref[IX.prop2.4]{2.4} (i), il existe une
   489	décomposition de $X$ en réunion finie de parties localement fermées
   490	constructibles $X_i$ telles que chaque $F|X_i$ soit localement
   491	constant et fini (\resp de présentation finie). Écrivons $X_i=U_i\cap
   492	\complement V_i$, où $U_i$ et $V_i$ sont des ouverts constructibles,
   493	et soit $N$ le nombre d'ouverts de $X$ dans la sous-topologie $T$
   494	engendrée par $\{U_i,V_i\}$. Raisonnons par récurrence sur $N$: si $W$
   495	est un élément non-vide minimal de $T$, il est évident que $F|W$
   496	est localement constant et de présentation finie. Soit $i:W\rightarrow
   497	X$ le morphisme d'inclusion et $Y=X-W$. On a la suite exacte
   498	$$
   499	0\longrightarrow i_{!}i^{\ast}F\longrightarrow F\longrightarrow F|
   500	Y\longrightarrow 0{\quoi},
   501	$$
   502	et on se réduit ainsi à le même assertion pour $Y$ et pour le faisceau
   503	$F| Y$, où on peut appliquer l'hypothèse de récurrence.
   504	
   505	\begin{enonce*}{Proposition 2.6}\etiquette[2.6]{IX.prop2.6}
   506	Supposons que $A$ soit un anneau noethérien.
   507	\begin{itemize}
   508	\item[(i)] Une limite projective (\resp inductive) finie de faisceaux de
   509	  $A$-modules ou d'ensembles constructibles est constructible. En
   510	  particulier, le noyau, le conoyau et l'image d'un morphisme de
   511	  faisceaux de $A$-modules constructibles sont constructibles.
   512	
   513	\item[(i bis)] Une limite projective finie de faisceaux d'ensembles
   514	  (\resp de groupes) constructibles est constructible, et si
   515	  $f:F\rightarrow G$ est un morphisme de faisceaux de groupes
   516	  constructibles, le noyau, l'image, et le conoyau (si $\Iim(f)$ est un
   517	  sous-groupe normal) sont constructibles. Une limite inductive finie
   518	  de faisceaux d'ensembles constructibles est constructible.
   519	
   520	\item[(ii)] Soit $0\rightarrow M'\rightarrow M\rightarrow M''\rightarrow 0$
   521	  une suite exacte de faisceaux de $A$-modules (\resp de groupes) avec
   522	  $M'$, $M''$ constructibles. Alors $M$ est constructible.
   523	
   524	\item[(iii)] Soit $M_1\rightarrow M_2\rightarrow M_3\rightarrow
   525	  M_4\rightarrow M_5$ une suite exacte de faisceaux de
   526	  $A$-modules. Alors si $M_i$ est constructible pour $i=1,2,4,5$, il
   527	  en est de même pour $i=3$.
   528	\end{itemize}
   529	\end{enonce*}
   530	
   531	\emph{Démonstration.} (i) Nous\pageoriginale laissons les assertions
   532	ensemblistes et pour les groupes au lecteur. Pour le cas des faisceaux
   533	de $A$-modules, il suffit de démontrer que le noyau et le conoyau d'un
   534	morphisme $f:F\rightarrow G$ de faisceaux de $A$-modules
   535	constructibles sont constructibles. L'assertion est locale sur $X$, et
   536	on peut donc supposer $X$ affine, donc quasi-compact et
   537	quasi-séparé. Alors on se réduit au cas où $F$ et $G$ sont localement
   538	constants par \hyperref[IX.prop2.4]{2.4}. (i), et on termine en utilisant \hyperref[IX.lem2.1]{2.1} (ii) et
   539	l'hypothèse noethérienne sur $A$. On prouve de façon analogue
   540	l'assertion (ii) en utilisant \hyperref[IX.lem2.1]{2.1} (iii). L'assertion (iii) est
   541	immédiate à partir de (i) et (ii).
   542	
   543	\begin{enonce*}{Proposition 2.7}\etiquette[2.7]{IX.prop2.7}
   544	Soient $X$ un schéma quasi-compact et quasi-séparé, $F$ un faisceau
   545	d'ensembles (\resp de $A$-modules). Pour que $F$ soit constructible,
   546	il faut et il suffit qu'il soit isomorphe au conoyau d'un couple de
   547	morphismes $H\rightrightarrows G$, où $H$ et $G$ sont des faisceaux
   548	d'ensembles représentables par des schémas étales de présentation
   549	finie sur $X$ (\resp que $F$ soit isomorphe au conoyau d'un
   550	homomorphisme $𝒜_{V,X}\rightarrow 𝒜_{U,X}$, où $U$ et $V$ sont
   551	deux schémas étales de présentation finie sur $X$).
   552	\end{enonce*}
   553	
   554	%% A_V [et non 𝒜] dans 2.9 \XXX
   555	
   556	La suffisance résulte facilement de \hyperref[IX.prop2.6]{2.6} (i bis).
   557	
   558	Supposons que $F$ soit un faisceau d'ensembles
   559	constructibles. Utilisant l'existence des sommes infinies dans le site
   560	étale sur $X$, on voit que l'on peut trouver un épimorphisme
   561	$G\rightarrow F$, avec $G$ représentable par un schéma étale sur $X$,
   562	que l'on peut de plus supposer somme de schémas affines $G_i$ $(i\in I)$,
   563	donc séparé sur $S$. Pour toute partie finie $J$ de l'ensemble
   564	d'indices $I$, soit $G_J$ la somme des $G_i$ pour $i\in J$, et $F_J$
   565	son image dans $F$. Comme $G_J$ et $F$ sont constructibles, il en est
   566	de même de $F_J$ (\hyperref[IX.prop2.6]{2.6} (i bis)), donc l'ensemble $X_J$ des $x\in X$
   567	tels que les fibres de $F$ et $F_J$ en un point géométrique de $X$ sur
   568	$x$ soient égales, est localement constructible \hyperref[IX.prop2.4]{2.4} (iv)). Comme la
   569	famille des $X_J$ est croissante et de réunion $X$, et que $X$ est
   570	quasi-compact, l'un des $X_J$ est égal à $X$ (\EGA IV 1.9.9). Cela
   571	montre, quitte à remplacer $G$ par $G_J$, que l'on peut supposer $G$
   572	affine, donc séparé sur $X$, et quasi-compact sur $X$ puisque $X$ est
   573	quasi-séparé (\EGA IV 1.2.4). Soit alors $H=\fprod{G}{G}{F}$. C'est un
   574	sous-faisceau de $\fprod{G}{G}{S}$ qui est représentable, donc est
   575	lui-même représentable (VIII 6.1). Comme $G$ est de plus constructible
   576	(\hyperref[IX.prop2.6]{2.6} (i bis)),\pageoriginale il résulte encore de l'argument précédent
   577	qu'il est quasi-compact, donc (étant séparé sur $X$) de présentation
   578	finie sur $X$. Cela prouve la première assertion de \hyperref[IX.prop2.7]{2.7}, et la
   579	deuxième se prouve par la même méthode (NB il n'est pas nécessaire que
   580	$A$ soit noethérien). Signalons en passant que la démonstration prouve
   581	aussi le corollaire suivant:
   582	
   583	\begin{enonce*}{Corollaire 2.7.1}\etiquette[2.7.1]{IX.cor2.7.1}
   584	Soient $X$ un schéma, $F$ un schéma étale sur $X$. Pour que le
   585	faisceau étale correspondant sur $X$ soit constructible, il faut et il
   586	suffit que $F$ soit de présentation finie sur $X$.
   587	\end{enonce*}
   588	
   589	La proposition \hyperref[IX.prop2.7]{2.7} nous sera surtout utile ici pour déduire certaines
   590	propriétés de passage à la limite pour les faisceaux constructibles:
   591	
   592	\begin{enonce*}{Corollaire 2.7.2}\etiquette[2.7.2]{IX.cor2.7.2}
   593	Soit $X$ un schéma quasi-compact et quasi-séparé. Alors tout faisceau
   594	d'ensembles (\resp de ind-p-groupes, \resp de $A$-modules) sur $X$ est
   595	limite inductive d'un système inductif filtrant de faisceaux
   596	d'ensembles (\resp\ldots) constructibles.
   597	\end{enonce*}
   598	
   599	Si $F$ est un faisceau d'ensembles, reprenons les notations de la
   600	démonstration de \hyperref[IX.prop2.7]{2.7}, et soit $L$ l'ensemble des couples $(J,H')$, où
   601	$J$ est une partie finie de $I$ et $H'$ une partie ouverte
   602	quasi-compacte de $H_J=\fprod{G_J}{G_J}{F}$. On ordonne $L$ de la fa\c
   603	con évidente, et on constate que, puisque $F$ est limite des $F_J$, et
   604	que chaque $F_J$ est limite des $\Coker(H'\rightarrow G_J)$ pour les
   605	ouverts quasi-compacts $H'$ de $H_J$, que $F$ est limite inductive du
   606	système inductif des $F_{\ell}=\Coker(H'\rightarrow F_J)$ pour
   607	$\ell=(H',J)\in L$, d'où le cas ensembliste. Le cas des faisceaux de
   608	$A$-modules se traite essentiellement de la même façon, en
   609	commençant par représenter $F$ comme conoyau d'un homomorphisme
   610	$𝒜_{V,X}\rightarrow 𝒜_{U,X}$, avec $U$, $V$ schémas étales et
   611	séparés sur $X$. Le cas des ind-$\sheaf{p}$-groupes demande un peu plus
   612	d'attention. Reprenant les notations précédentes, nous désignons par
   613	$L(G_J)$ le faisceau en groupes libre engendré par le faisceau
   614	d'ensembles $G_J$, par $N_J$ le noyau de l'homomorphisme
   615	$L(G_J)\rightarrow F$ déduit de $G_J\rightarrow F$, par $L$ l'ensemble
   616	des couples $(J,N')$, où $N'$ est un sous-faisceau d'\emph{ensembles}
   617	de $N_J$ qui est « quasi-compact » \ie tel qu'il existe un
   618	épimorphisme $P\rightarrow N'$, avec $P$ représentable par un
   619	$X$-schéma quasi-compact. On ordonne $L$ en déclarant que $(J,N')$ est
   620	plus petit que $(J_1, N'_1)$ si $J\subset J_1$ et si
   621	l'homomorphisme\pageoriginale canonique $L(G_J)\rightarrow L(G_{J_1})$
   622	applique $N_1$ dans $N'_1$. Il est immédiat que $L$ est filtrant, et
   623	que l'on obtient un système inductif de faisceaux en groupes
   624	$F_{\ell}$, indexé par $L$, en prenant, pour $\ell=(J,N')$,
   625	$F_{\ell}=$ faisceau en groupes quotient de $L(G_J)$ par le
   626	sous-faisceau en groupes invariant engendré par $N'$. Il est immédiat
   627	de plus que $F$ est isomorphe à la limite inductive des $F_{\ell}$, et
   628	il reste à prouver qu'il existe une partie $L'$ de $L$, cofinale dans
   629	$L$, telle que pour $\ell\in L'$, $F_{\ell}$ soit un faisceau en
   630	$\sheaf{p}$-groupes \emph{constructible}. Pour ceci, il suffit de prouver
   631	que pour tout $J$, on peut trouver un sous-faisceau d'ensembles $N'$
   632	de $N_J$ qui soit quasi-compact, contienne un sous-faisceau
   633	d'ensembles quasi-compact donné $N'_0$, et soit tel que le $F_{\ell}$
   634	correspondant soit un faisceau en $\sheaf{p}$-groupes
   635	constructible. Notons d'abord qu'il suffit, pour ceci, que les fibres
   636	de $F_{\ell}$ soient des $\sheaf{p}$-groupes finis: alors $F_{\ell}$ sera
   637	automatiquement constructible, comme le lecteur vérifiera sans
   638	peine. D'autre part, le lecteur pourra vérifier également que
   639	l'ensemble $X_{N'}$ des points $x$ de $X$ tels que la fibre de
   640	$F_{\ell}$ en un point géométrique au-dessus de $x$ soit un
   641	$\sheaf{p}$-groupe fini est localement constructible. De plus, les parties
   642	$X_{N'}$ de $X$ sont fonction croissante de $N'$, enfin leur réunion
   643	est $X$ tout entier: ce dernier point résulte du fait que la limite
   644	inductive des $F_{\ell}$, pour $\ell=(N',J)$ avec $J$ fixé, est le
   645	sous-faisceau en groupes $F_J$ de $F$ engendré par le sous-faisceau
   646	d'ensembles $F_J$, qui est quasi-compact, donc $\bar{F_J}$ est à
   647	fibres des $\sheaf{p}$-groupes finis. D'autre part, dans un groupe libre à
   648	un nombre fini de générateurs, donc s'il est représenté comme
   649	réunion filtrante de ses parties finies, l'une de celles-ci engendre
   650	tout le sous-groupe. De là résulte aisément que la réunion des
   651	$X_{N'}$ est bien $X$. Utilisant à nouveau (\EGA IV 1.9.9) on trouve
   652	que un des $X_{N'}$ est égal à $X$, ce qui achève la démonstration de
   653	\hyperref[IX.cor2.7.2]{2.7.2}.
   654	
   655	\setcounter{subsection}{7}
   656	\setcounter{subsubsection}{2}
   657	\paragraph{}\etiquette[2.7.2.1]{IX.par2.7.2.1}
   658	Notons que la démonstration un peu pénible du cas des schémas en
   659	groupes met en évidence le manque d'un sorite commode pour les
   660	conditions de constructibilité pour des faisceaux en groupes à fibres
   661	pas nécessairement finies; nous nous en excusons auprès du lecteur en
   662	l'invitant à combler au besoin cette lacune par ses propres moyens, et
   663	lui signalons en guise de consolation que dans \hyperref[IX.prop2.9]{2.9} (iii) nous obtenons
   664	une démonstration plus simple lorsque $X$ est noethérien.
   665	
   666	\begin{enonce*}{Corollaire 2.7.3}\etiquette[2.7.3]{IX.cor2.7.3}
   667	Soient\pageoriginale $X$ un schéma quasi-compact et quasi-séparé, $F$
   668	un faisceau d'ensembles (\resp de ind-groupes finis, \resp de
   669	$A$-modules) constructible. Alors le foncteur $\Hom(F,G)$, où $G$ est
   670	un faisceau d'ensembles (\resp\ldots) variable, commute aux limites
   671	inductives en $G$. Dans le cas où $F$ est un faisceau de $A$-modules,
   672	$A$ étant noethérien, les foncteurs $\Ext^{i}(X;F,G)$ en le $A$-module
   673	$G$ commutent également aux limites inductives.
   674	\end{enonce*}
   675	
   676	Ceci résulte de \hyperref[IX.prop2.7]{2.7} et (VII 3.3) par les arguments habituels, qui sont
   677	laissés au lecteur. De même, la conjonction de \hyperref[IX.prop2.7]{2.7} et des résultats de
   678	(VII 5) donne aisément:
   679	
   680	\begin{enonce*}{Corollaire 2.7.4}\etiquette[2.7.4]{IX.cor2.7.4}
   681	Soit $I$ une catégorie filtrante, $i\mapsto X_i:I^0\rightarrow (\Sch)$
   682	un foncteur qui transforme les flèches en morphismes affines et les
   683	objets en des schémas quasi-compacts et quasi-séparés. Pour tout
   684	schéma $Y$, soit $F_{c}(Y)$ la catégorie des faisceaux d'ensembles
   685	(\resp de ind-$\sheaf{p}$-groupes, \resp de $A$-modules) constructibles
   686	sur $Y$. Alors on a une équivalence de catégories
   687	$$
   688	F_{c}(X)\xleftarrow{\approx}\varinjlim_i F_{c}(X_i){\quoi},
   689	$$
   690	où $X=\varprojlim_i X_i$ {\rm (VII 5)} et où le deuxième membre
   691	désigne la catégorie $\varinjlim$ de catégories fibrées, définie dans\hspace{1cm}
   692	En particulier, tout faisceau d'ensembles (\resp\ldots) constructible
   693	sur $X$ est isomorphe à l'image inverse d'un faisceau de même nature
   694	sur un des $X_i$.
   695	\end{enonce*}
   696	
   697	\begin{enonce*}{Proposition 2.8}\etiquette[2.8]{IX.prop2.8}
   698	Soit $f:X\rightarrow Y$ un morphisme surjectif et localement de
   699	présentation finie, et soit $F$ un faisceau d'ensembles (\resp\ldots)
   700	sur $Y$. Alors $F$ est constructible si et seulement si $f^{\ast}(F)$ l'est.
   701	\end{enonce*}
   702	
   703	Nous utiliserons le
   704	
   705	\begin{enonce*}{Lemme 2.8.1}\etiquette[2.8.1]{IX.lem2.8.1}
   706	Soit $f:X\rightarrow Y$ un morphisme surjectif et localement de
   707	présentation finie, avec $Y$ quasi-compact et quasi-séparé. Alors il
   708	existe une partition finie de $Y$ en des sous-schémas $Y_i$ de
   709	présentation finie sur $Y$ (en particulier\pageoriginale chaque $Y_i$
   710	est constructible dans $X$) et pour tout $i$, des morphismes finis
   711	surjectifs $Y''_i\xrightarrow{g_i}Y'_i\xrightarrow{h_i}Y_i$, avec
   712	$h_i$ étale et $g_i$ localement libre (\ie plat et de présentation
   713	finie, en plus de la condition $g_i$ fini), et radiciel, et enfin un
   714	$Y$-morphisme $Y''_i\rightarrow X$\footnote[1]{Cet énoncé est aussi
   715	  donné dans \EGA IV 17.16.4.}.
   716	\end{enonce*}
   717	
   718	Utilisant le fait que $Y$ est quasi-compact et quasi-séparé, on se
   719	ramène aisément au cas où $Y$ est quasi-affine donc séparé (en
   720	utilisant un recouvrement ouvert affine fini $(U_i)_{1\leq i\leq n}$
   721	de $Y$ et considérant les $Y_i=U_i-U_i\cap \bigcup_{j< i} U_j$), puis
   722	au cas $Y$ affine (par le même argument). Le schéma $X$ est réunion
   723	d'ouverts affines $X_i$, et l'image $T_i$ de $X_i$ dans $Y$ est
   724	constructible (\EGA IV 1.8.4), la réunion des $T_i$ est $X$ puisque $f$
   725	est surjectif, d'où s'ensuit que $X$ est réunion d'un nombre fini des
   726	$T_i$ (\EGA IV 1.9.9), de sorte que, remplaçant $X$ par le schéma
   727	somme des $X_i$ correspondants, on peut supposer $X$ affine. Alors le
   728	procédé de passage à la limite standard (\EGA IV 8 8.1.2 c)) nous
   729	permet de supposer de plus $Y$ noethérien. La récurrence noethérienne
   730	habituelle, et le procédé de passage à la limite (\EGA IV 8.1.2 a))
   731	nous ramène au cas où $Y$ est réduit au spectre d'un corps $k$. Comme
   732	alors $X\neq ∅$, il existe un point fermé $x$ dans $X$, qui
   733	correspond à une extension finie $k'$ de $k$, elle-même extension
   734	radicielle d'une extension séparable $k'_s$ de $k$. On prendra alors
   735	$Y'=\Spec(k'_s)$, $Y''=\Spec(k')$, ce qui achève la démonstration de
   736	\hyperref[IX.lem2.8.1]{2.8.1}.
   737	
   738	La nécessité dans \hyperref[IX.prop2.8]{2.8} étant triviale (\hyperref[IX.prop2.4]{2.4} (iii)), prouvons la
   739	suffisance. On est alors ramené, par le lemme précédent, au cas où $f$
   740	est de la forme $gh$, où $g$ et $h$ satisfont aux conditions énoncées
   741	pour $g_i$, $h_i$ dans le lemme. On peut donc supposer successivement,
   742	soit que $f$ est fini radiciel, cas trivial par (VIII 1.1), soit que
   743	$f$ est fini étale. Notons d'ailleurs que, utilisant l'hypothèse que
   744	$f^{\ast}(F)$ est constructible et la définition de la
   745	constructibilité, nous pouvions supposer (quitte à remplacer $X$ par
   746	un $X$-préschéma $X'=\coprod_i X_i$ de présentation finie et à
   747	morphisme structural surjectif) que $f^{\ast}(F)$ est déjà localement
   748	constant. Mais lorsque $f$ est étale et surjectif, cela implique que
   749	$F$ est localement constant. Comme de plus les hypothèses faites sur
   750	les fibres de $f^{\ast}(F)$ restent valables pour celles de $F$, $f$
   751	étant surjectif, il s'ensuit bien que $F$ est constructible, ce qui
   752	achève la démonstration de \hyperref[IX.prop2.8]{2.8}.
   753	
   754	\begin{enonce*}{Proposition 2.9}\etiquette[2.9]{IX.prop2.9}
   755	Soient\pageoriginale $X$ un schéma noethérien, $A$ un anneau
   756	noethérien, et $p$ un ensemble de nombres premiers.
   757	\begin{enumerate}
   758	\renewcommand{\theenumi}{\roman{enumi}}
   759	\renewcommand{\labelenumi}{(\theenumi)}
   760	\item La catégorie des faisceaux d'ensembles (\resp de ind
   761	$p$-groupes, \resp de $A$-modules) sur $X$ est localement
   762	noethérienne, c'est-à-dire, possède un ensemble de générateurs formé
   763	d'objets noethériens.
   764	
   765	\item Un faisceau $F$ d'ensembles (\resp de ind $p$-groupes, \resp de
   766	$A$-modules) sur $X$ est constructible si et seulement si il est
   767	noethérien. Si $F$ est un faisceau de $A$-modules, alors $F$ est
   768	constructible si et seulement si il est quotient d'une somme finie de
   769	faisceaux de la forme $A_{U/X}$, où $U\rightarrow X$ est étale et de
   770	type fini (notation de {\rm IV.2.4}: $A_{U/X}=A_U$).
   771	
   772	%% vérifier la référence IV.2.4 [me semble incorrecte] \XXX
   773	
   774	\item Les sous-faisceaux constructibles d'un faisceau $F$ d'ensembles
   775	(\resp\ldots) forment un système inductif, et $F$ est limite inductive
   776	de ces sous-faisceaux.
   777	\end{enumerate}
   778	\end{enonce*}
   779	
   780	Démonstration. On laisse l'assertion ensembliste au lecteur. Traitons
   781	d'abord le cas d'un faisceau de $A$-modules. Pour (i) il faut trouver
   782	un ensemble de générateurs dont les éléments sont des objets
   783	noethériens. Or les faisceaux de la forme $A_{U/X}$, $U\rightarrow X$
   784	étale et de type fini forment un ensemble de générateurs dont on voit
   785	facilement qu'ils sont constructibles, et il suffit donc de démontrer
   786	le lemme suivant:
   787	
   788	\begin{enonce*}{Lemme 2.10}\etiquette[2.10]{IX.lem2.10}
   789	Soit $X$ noethérien. Alors tout $A$-Module constructible est noethérien.
   790	\end{enonce*}
   791	
   792	Démonstration du lemme.
   793	
   794	Par récurrence noethérienne nous supposons le lemme vrai pour chaque
   795	sous-schéma $X'$ de $X$ distinct de $X$. Prenons un ouvert non-vide
   796	$X_0$ de $X$ tel que $F$ induise un faisceau localement constant $F_0$
   797	sur $X_0$.
   798	
   799	Soit maintenant $\{F_i\}$, $i\in \NN$ une suite
   800	croissante de sous-faisceaux de $F_0$. Comme pour un point géométrique
   801	$P$, générique pour une composante irréductible de $X_0$, la fibre de
   802	$F_0$ en $P$ est de type fini, la suite des fibres
   803	$(F_i)_P$\pageoriginale est stationnaire, et nous pouvons supposer
   804	qu'elle est constante. Or soit $Q$ un point géométrique de $X_0$
   805	\emph{spécialisation} de $P$ (VIII 7.2). Puisque $F_0$ est localement
   806	constant sur $X_0$, on voit immédiatement que le morphisme de
   807	spécialisation $F_{0_Q}\rightarrow F_{0_P}$ (VIII.7) est bijectif,
   808	donc $(F_i)_Q\rightarrow (F_i)_P$ est \emph{injectif} pour chaque $i$.
   809	
   810	Soient $V/X$ étale et $P$-ponctué, et $s_1,\ldots,s_n$ des sections
   811	$\in F_1(V)$ qui engendrent $(F_1)_P$. Alors les $s_i$ engendrent
   812	$(F_i)_Q$ pour chaque $i$ et chaque point géométrique $Q$
   813	spécialisation de $P$ au-dessus de $V$. La suite de faisceaux $F_i$
   814	est donc constante dans un voisinage de l'image de $P$, disons dans
   815	$X'\neq ∅$. Soit $Y=X-X'$, il suffit donc de prouver que la suite
   816	des $F_i| Y$ est stationnaire, ce qui résulte de l'hypothèse de
   817	récurrence noethérienne.
   818	
   819	Les assertions (ii) pour les $A$-modules sont maintenant immédiates:
   820	Comme les $A_{U/X}$ sont des générateurs noethériens, on a évidemment
   821	$F$ noethérien $\Leftrightarrow F$ est quotient d'une somme finie de
   822	faisceaux $A_{U/X}$. Or $A_{U/X}$ est constructible, et par suite,
   823	compte tenu de \hyperref[IX.prop2.6]{2.6}, $F$ noethérien $\Rightarrow F$ constructible;
   824	l'implication inverse étant déjà établie (\hyperref[IX.lem2.10]{2.10}), cela établit
   825	(ii). Enfin, l'assertion (iii) pour les $A$-modules est une propriété
   826	des catégories localement noethériennes (\cf [1], Ch. IV).
   827	
   828	Traitons maintenant le cas des ind-$p$-groupes. Notons d'abord qu'un
   829	faisceau de groupes constructible $F$ est certainement noethérien. En
   830	effet, puisque les fibres sont finies, on peut employer le même
   831	argument qu'avec les $A$-modules. On va maintenant démontrer
   832	l'assertion (iii) directement, ce qui impliquera (i) - en effet, cela
   833	démontrera que les faisceaux constructibles forment un ensemble de
   834	générateurs. De plus, un faisceau noethérien quelconque sera alors
   835	nécessairement constructible, d'où (ii).
   836	
   837	Pour démontrer (iii), il suffit évidemment de démontrer qu'un
   838	sous-faisceau (en groupes) $S$ d'un faisceau $F$ de ind $p$-groupes
   839	engendré par un nombre fini de sections $s_i\in F(U_i)$,
   840	$U_i\rightarrow X$ étale de type fini, est constructible
   841	\footnote[1]{Rappelons la définition de ce faisceau: soit $L_{U/X}$
   842	le faisceau qui représente le foncteur $F\mapsto F(U)$ dans la
   843	catégorie des groupes. La section $s_i\in F(U_i)$ induit un morphisme
   844	$L_{U_i/X}\rightarrow F$, et $S$ est l'image du morphisme somme
   845	$\amalg_i L_{U_i/X}\rightarrow F$. C'est le plus petit sous-faisceau
   846	qui « contient » les sections $s_i$.}. (Puisqu'un\pageoriginale
   847	faisceau constructible est noethérien, il est engendré par un nombre
   848	fini de sections.) Par récurrence noethérienne, on peut supposer que
   849	c'est vrai pour chaque sous-schéma fermé $Y$ de $X$, distinct de
   850	$X$. Or la notion de sous-faisceau engendré par des sections commute
   851	avec l'image inverse de faisceaux; en effet, soit $L_{U/X}$ le
   852	faisceau qui représente le foncteur $F\mapsto F(U)$ dans la catégorie
   853	des groupes. Alors si $f:Y\rightarrow X$, on a
   854	$f^{\ast}(L_{U/X})=L_{\fprod{L}{U}{X}/Y}$. Comme $S$ est l'image de la
   855	somme des $L_{U_i/X}$, et comme $f^{\ast}$ commute aux sommes, il est
   856	clair que $f^{\ast}S$ est le sous-faisceau de $f^{\ast}F$ engendré par
   857	les sections $f^{\ast}(S_i)$.
   858	
   859	On peut donc supposer que $S$ induit un faisceau constructible sur
   860	chaque fermé $Y$ distinct de $X$, et il suffit ainsi (\hyperref[IX.prop2.4]{2.4} (i)) de
   861	démontrer que $S$ induit un faisceau constructible sur un ouvert
   862	non-vide convenable de $X$. En remplaçant $X$ par un ouvert assez
   863	petit, on peut supposer que chaque $U_i$ est un revêtement étale de
   864	$X$. Il suffit de démontrer qu'alors $S$ est localement constant, à
   865	fibres finies sur un ouvert non vide convenable. C'est une assertion
   866	locale sur $X$ pour la topologie étale, et on peut donc supposer que
   867	chaque $U_i$ est complètement décomposé, disons $U_i=\amalg_{n_i}X$
   868	(somme de $n_i$ copies de $X$). Soient $s_{ij}$ $(j=1,\ldots,n_i)$ les
   869	sections de $F(X)$ telles que
   870	$s_i$ $(s_{il},\ldots,s_{in_i})$. Évidemment, $S$ n'est autre que le
   871	sous-faisceau de $F$ engendré par les $s_{ij}\in F(X)$. Comme $F(X)$
   872	est un ind $p$-groupe (\hyperref[IX.def1.5]{1.5}) les $s_{ij}$ engendrent un sous-groupe
   873	fini \hyperref[IX.def1.3]{1.3}. Donc on peut supposer que l'ensemble des $s_i$ est un
   874	sous-groupe fini de $F(X)$, mais alors $S$ est identique au faisceau
   875	\emph{d'ensembles} engendré par les $s_i$, qui est constructible en
   876	vertu de \hyperref[IX.prop2.6]{2.6} (i) comme faisceau d'ensembles, donc comme faisceau de
   877	groupes, \cqfd
   878	
   879	Dans les deux propositions suivantes, nous donnons des critères pour
   880	qu'un faisceau soit localement constant, ou constructible, en
   881	utilisant les homomorphismes de spécialisation (VIII 7.7).
   882	
   883	\begin{enonce*}{Proposition 2.11}\etiquette[2.11]{IX.prop2.11}
   884	Soient $X$ un schéma, $F$ un faisceau d'ensembles (\resp de groupes
   885	ind-finis, \resp de $A$-modules) constructible; dans le cas
   886	ensembliste,\pageoriginale on suppose de plus $F$ à fibres
   887	finies. Soient $x\in X$, $x$ un point géométrique sur $x$. Pour que
   888	$F$ soit localement constant au voisinage de $x$, il faut et il suffit
   889	que pour tout point géométrique $\bar{x}'$ générisation de $\bar{x}$,
   890	et toute flèche de spécialisation $\bar{x'}\rightarrow \bar{x}$ {\rm
   891	(VIII 7.2)}, l'homomorphisme de spécialisation {\rm (VIII 7.7)}
   892	$F_{\bar{x}}\rightarrow F_{\bar{x'}}$ soit un isomorphisme.
   893	\end{enonce*}
   894	
   895	Le cas d'un faisceau en groupes étant un cas particulier de celui d'un
   896	faisceau d'ensembles, nous nous contentons de traiter celui-ci; le cas
   897	d'un faisceau de modules se traite de façon essentiellement
   898	identique. La nécessité de la condition étant triviale, nous nous
   899	bornons à établir la suffisance. Comme la fibre $I=F_{\bar{x}}$ est
   900	finie, quitte à remplacer $X$ par un schéma $X'$ étale sur $X$
   901	convenable, on peut supposer que l'on peut trouver un homomorphisme
   902	$u:I_X\rightarrow F$ du faisceau constant $I_X$ dans $F$, induisant un
   903	isomorphisme pour les fibres en $\bar{x}$. Utilisant la
   904	constructibilité des deux faisceaux, et \hyperref[IX.prop2.6]{2.6}, on voit aisément que
   905	l'ensemble $Z$ des points $z$ de $X$ tels que l'homomorphisme induit
   906	sur les fibres en un point géométrique $\bar{z}$ sur $z$ soit
   907	bijective, est une partie localement constructible de $X$. L'hypothèse
   908	implique aussitôt qu'elle contient les générisations de $x$, donc (\EGA
   909	IV 1.10.1) c'est un voisinage de $x$, ce qui prouve que $F$ est
   910	localement constant (en l'occurrence, même constant) au voisinage de
   911	$x$,\cqfd
   912	
   913	\begin{enonce*}{Définition 2.12}\etiquette[2.12]{IX.def2.12}
   914	\footnote[1]{La terminologie est due à S. Lubkin.}: Un schéma $X$
   915	est dit connexe par arcs si pour tout couple $P$, $Q$ de points
   916	géométriques de $X$, il existe des points géométriques
   917	$P=P_0,\ldots,P_n$; $Q_1,\ldots,Q_n=Q$ et des spécialisations {\rm
   918	(VIII 7.2)} $P_i\rightarrow Q_i$ et $P_{i-1}\rightarrow
   919	Q_i$ $(i=1,\ldots,n)$. On dit que $X$ est localement connexe par arcs
   920	(pour la topologie étale) si pour chaque $U\rightarrow X$ étale il
   921	existe un recouvrement étale $\{U_i\rightarrow U\}$ de $U$ tel que
   922	$U_i$ est connexe par arcs pour chaque $i$.
   923	\end{enonce*}
   924	
   925	On vérifie immédiatement qu'on préschéma localement noethérien est
   926	localement connexe par arcs.
   927	
   928	\begin{enonce*}{Proposition 2.13}\etiquette[2.13]{IX.prop2.13}
   929	\begin{enumerate}
   930	\renewcommand{\theenumi}{\roman{enumi}}
   931	\renewcommand{\labelenumi}{(\theenumi)}
   932	\item Soit\pageoriginale $X$ localement connexe par arcs et soit $F$
   933	un faisceau d'ensembles (\resp de groupes, \resp de $A$-modules) sur
   934	$X$. Supposons que les fibres de $F$ sont finies (\resp finies, \resp
   935	de présentation finie). Alors $F$ est localement constant si et
   936	seulement si pour toute spécialisation $P\rightarrow Q$ de points
   937	géométriques, le morphisme de spécialisation {\rm (VIII 7.7)}
   938	$F_Q\rightarrow F_P$ est bijectif.
   939	
   940	\item Supposons que $X$ soit localement noethérien et que les fibres
   941	de $F$ soient finies (\resp finies, \resp de présentation
   942	finie). Alors, si pour toute spécialisation $P\rightarrow Q$, le
   943	morphisme de spécialisation $F_Q\rightarrow F_P$ est injectif, $F$ est
   944	constructible.
   945	
   946	\item Supposons que $X$ soit localement noethérien et que les fibres
   947	du faisceau d'ensembles $F$ soient finies. Soit $c:X\rightarrow
   948	\NN$ la fonction qui à $x\in X$ associe le nombre
   949	d'éléments de la fibre $F_{\bar{x}}$ de $F$ en un point géométrique
   950	$\bar{x}$ au-dessus de $x$. Alors $F$ est constructible si et seulement
   951	si $c$ est une fonction constructible \ie, si et seulement si
   952	$c^{-1}(n)$ est constructible pour chaque $n\in \NN$.
   953	\end{enumerate}
   954	\end{enonce*}
   955	
   956	\emph{Démonstration.}
   957	\begin{enumerate}
   958	\renewcommand{\theenumi}{\roman{enumi}}
   959	\renewcommand{\labelenumi}{(\theenumi)}
   960	\item Soit $P$ un point géométrique de $X$. On va trouver
   961	un voisinage étale de $P$, c'est-à-dire, un $U\rightarrow X$ étale
   962	$P$-ponctué, tel que $F$ devienne constant sur $U$. Soit $V'$ un
   963	voisinage étale de $P$ tel que $F(V)\rightarrow F_P$ soit
   964	surjectif. Un tel $V'$ existe d'après l'hypothèse de finitude. Soit
   965	$U\rightarrow V'$ un voisinage étale de $P$ dans $V'$, tel que $U$
   966	soit connexe par arcs. Alors $F(U)\rightarrow F_P$ est encore
   967	surjectif. Choisissons un sous-ensemble $S\subset F(U)$ qui s'envoie
   968	bijectivement sur $F_P$. Si $F$ est un faisceau de groupes (\resp de
   969	$A$-modules), on voit facilement qu'on peut, en changeant au besoin
   970	$U$, trouver un tel ensemble qui soit de plus un sous-groupe (\resp
   971	sous-module) de $F(U)$. On déduit un morphisme $S_U\rightarrow F|
   972	U$, où $S_U$ est le faisceau constant à valeur $S$. On a
   973	$$
   974	(S_U)_P\xrightarrow{\sim}F_P{\quoi}.
   975	$$
   976	Puisque\pageoriginale chaque morphisme de spécialisation sur les
   977	fibres de $F$ (\resp $S_U$) est bijectif, et puisque $U$ est connexe
   978	par arc, il s'ensuit que $S_U\xrightarrow{\sim}F| U$, d'où le
   979	résultat.
   980	
   981	\item L'assertion est locale, et on peut donc supposer $X$
   982	noethérien. Soit $P$ un point géométrique au-dessus d'un point maximal
   983	de $X$. Il existe un voisinage étale irréductible, donc connexe par
   984	arcs, $U\rightarrow X$ de $P$ tel que $F(U)$ s'envoie surjectivement
   985	sur $F_P$. Il s'envoie alors surjectivement sur $F_Q$ pour chaque
   986	point géométrique $Q$ de $V$, ou ce qui revient au même, pour chaque
   987	point géométrique de l'image $U$ de $V$ dans $X$, car un tel point est
   988	spécialisation de $P$ et $F_Q\rightarrow F_P$ est injectif. Donc
   989	chaque morphisme de spécialisation dans $U$ est bijectif, donc $F|
   990	U$ est localement constant. Par récurrence noethérienne on peut
   991	supposer que $F| X-U$ est constructible, d'où le résultat.
   992	
   993	\item Il suffit évidemment de traiter le cas où la fonction $c$ est
   994	constante. Soit $V\rightarrow X$ un voisinage étale d'un point
   995	géométrique $P$ au-dessus d'un point maximal de $X$, tel que $F(V)$
   996	s'envoie surjectivement sur $F_P$. Soit $S$ un sous-ensemble de $F(V)$
   997	tel que $S\xrightarrow{\sim}F_P$. Alors on a, pour le faisceau
   998	constant $S_V$ et pour chaque spécialisation $P\rightarrow Q$ de
   999	points géométriques, un diagramme commutatif
  1000	$$
  1001	\xymatrix@C=2cm@R=1.25cm{
  1002	(S_V)_P   \ar[r]^{\sim}& F_P\\
  1003	(S_V)_Q   \ar[r] \ar[u]^{\text{\rotatebox{90}{$\backsim$}}}  & F_Q. \ar[u]
  1004	}
  1005	$$
  1006	Comme $c(F_Q)=c(F_P)$, la flèche $F_Q\rightarrow F_P$ est bijective,
  1007	donc $F$ est localement constant dans un voisinage de $P$ d'après
  1008	(i). Par récurrence noethérienne, on a gagné.
  1009	\end{enumerate}
  1010	
  1011	Le résultat (ii) suivant aura une utilité technique dans la suite; il
  1012	permettra de réduire certaines vérifications au cas d'un faisceau
  1013	constant.
  1014	
  1015	\begin{enonce*}{Proposition 2.14}\etiquette[2.14]{IX.prop2.14}
  1016	Soit\pageoriginale $X$ un schéma.
  1017	\begin{enumerate}
  1018	\renewcommand{\theenumi}{\roman{enumi}}
  1019	\renewcommand{\labelenumi}{(\theenumi)}
  1020	\item Pour tout morphisme fini et de présentation finie
  1021	$f:X'\rightarrow X$, et tout faisceau d'ensembles (\resp de groupes
  1022	ind-finis, \resp de $A$-modules) constructible sur $X'$, $f_{\ast}(F)$
  1023	est constructible.
  1024	
  1025	\item Supposons $X$ quasi-compact et quasi-séparé, et soit $F$ un
  1026	faisceau d'ensembles (\resp\ldots) constructible sur $X$. Alors on
  1027	peut trouver une famille finie $(p_i:X'_i\rightarrow X_i)$ de
  1028	morphismes finis, pour tout $i$, un faisceau d'ensembles (\resp\ldots)
  1029	constructible constant $C_i$ sur $X'_i$, et un monomorphisme
  1030	$$
  1031	F\hookrightarrow \prod_i p_{i\ast}(C_i){\quoi}.
  1032	$$
  1033	\end{enumerate}
  1034	\end{enonce*}
  1035	
  1036	Prouvons d'abord (i). Par définition, on peut supposer $X$ affine, et
  1037	utilisant alors \hyperref[IX.cor2.7.4]{2.7.4}, le procédé de passage à la limite standard nous
  1038	ramène au cas où $X$ est noethérien. Un nouveau passage à la limite,
  1039	utilisant encore \hyperref[IX.cor2.7.4]{2.7.4} et de plus \hyperref[IX.prop2.4]{2.4}. (v), nous ramène au cas où $X$ est le spectre d'un corps. Alors la conclusion se réduit à l'assertion
  1040	correspondante sur la fibre de $f_{\ast}(F)$ en un point géométrique
  1041	sur $X$, laquelle résulte immédiatement de (VIII 5.5, 5.8).
  1042	
  1043	Prouvons (ii). En vertu de \hyperref[IX.prop2.4]{2.4} (i), il existe une partition de $X$ en
  1044	sous-schémas $Y_i$, avec des morphismes d'inclusion
  1045	$u_i:Y_i\rightarrow X$ de présentation finie, tels que la restriction
  1046	$F_i$ de $F$ à chaque $Y_i$ soit un faisceau localement constant. Il
  1047	est alors immédiat que les homomorphismes canoniques $F\rightarrow
  1048	u_{i\ast}(F_i)$ définissent un homomorphisme $F\rightarrow \prod_i
  1049	u_{i\ast}(F_i)$ qui est un monomorphisme.
  1050	
  1051	On pouvait même s'arranger dans la construction précédente pour que
  1052	chaque $F_i$ soit, non seulement localement constant, mais de plus
  1053	\emph{isotrivial,} \ie qu'il existe un morphisme étale \emph{fini}
  1054	surjectif $q_i:Y'_i\rightarrow Y_i$ tel que $q^{\ast}_i(F_i)$ soit un
  1055	faisceau constant sur $Y'_i$, soit $C_{iY'_i}$. Cela résulte par
  1056	exemple aisément de la définition de \hyperref[IX.lem2.8.1]{2.8.1} et de (VIII 1.1). D'autre
  1057	part, il\pageoriginale résulte du « Main Theorem »  sous la forme
  1058	(\EGA IV 8.12.6) qu'il existe un morphisme fini $p_i:Z_i\rightarrow X$
  1059	et une immersion ouverte $v_i:Y'_i\rightarrow Z_i$ de $X$-schémas.
  1060	
  1061	Comme $F_i$ se plonge dans
  1062	$q_{i\ast}(q^{\ast}_i(F_i))=q_{i\ast}(C_{iY'_i})$, $F$ se plonge dans
  1063	le produit des
  1064	$u_{i\ast}(q_{i\ast}(C_{iY'_i}))=p_{i\ast}(v_{i\ast}(C_{iY'_i}))$. Si
  1065	tout $Z_i$ est normal, alors le lemme (\hyperref[IX.lem2.14.1]{2.14.1}) ci-dessous implique que
  1066	$v_{i\ast}(C_{iY'_i})\simeq C_{iZ_i}$, donc $F$ se plonge dans le
  1067	produit des $p_{i\ast}(C_{iZ_i})$, et on a terminé. Dans le cas
  1068	général, introduisons le normalisé $\bar{Z}_i$ de $Z_i$, l'image
  1069	inverse $\bar{Y}'_i$ de $Y'_i$ dans $\bar{Z}_i$, l'immersion
  1070	$\bar{v}_i:\bar{Y}'_i\rightarrow \bar{Z}_i$. On voit aussitôt que
  1071	$\bar{Y}'_i$ contient les points maximaux de $\bar{Z}_i$, donc est
  1072	dense dans $\bar{Z}_i$. Si $r_i:\bar{Z}_i\rightarrow Z_i$ et
  1073	$s_i:\bar{Y}'_i\rightarrow Y'_i$ sont les projections, comme cette
  1074	dernière est surjective, il s'ensuit que $C_{iY'_i}$ se plonge dans
  1075	$s_{i\ast}(C_{i\bar{Y}'_i})$, dons $v_{i\ast}(C_{iY'_i})$ se plonge
  1076	dans
  1077	$v_{i\ast}(s_{i\ast}(C_{i\bar{Y}'_i}))=r_{i\ast}(\bar{v}_{i\ast}(C_{i\bar{Y}'_i}))$,
  1078	qui est isomorphe à $r_{i\ast}(C_{i\bar{Z}_i})$ en vertu de (\hyperref[IX.lem2.14.1]{2.14.1})
  1079	ci-dessous, donc $F$ se plonge dans le produit des
  1080	$p_{i\ast}(r_{i\ast}(C_{\bar{Z}_i}))=t_{i\ast}(C_{i\bar{Z}_i})$, où
  1081	$t_i=p_ir_i$ est la projection canonique $\bar{Z}_i\rightarrow
  1082	X$. Lorsque ce morphisme est fini pour tout $i$, on a donc
  1083	terminé. Dans le cas général, on peut dire seulement que $\bar{Z}_i$
  1084	est \emph{entier} sur $X$, donc de la forme $\sheaf{\Spec}(𝒜_i)$, où
  1085	$𝒜_i$ est un faisceau quasi-cohérent de $𝒪_X$-Algèbres, qui est
  1086	entier sur $𝒪_X$. Utilisant (\EGA IV 1.7.7), on voit que $𝒜_i$ est
  1087	limite inductive de ses sous-Algèbres $𝒜_{i,j}$ finies sur $X$, donc
  1088	$\bar{Z}_i$ est limite projective (VII 5) des préschémas
  1089	$Z_{ij}=\Spec(A_{ij})$, finis sur $X$. Appliquant (VII 5.11), on
  1090	trouve que $t_{i\ast}(C_{i\bar{Z}_i})$ est limite inductive des
  1091	$t_{ij^{\ast}}(C_{iZ_{ij}})$ où $t_{ij}:Z_{ij}\rightarrow X$ est la
  1092	projection canonique. En vertu de (\hyperref[IX.cor2.7.3]{2.7.3}), l'homomorphisme
  1093	$F\rightarrow t_{i\ast}(C_{i\bar{Z}_i})$ se factorise donc par un
  1094	homomorphisme $F\rightarrow t_{ij\ast}(C_{iZ_{ij}})$, pour un $j=j(i)$
  1095	convenable. On pose maintenant $X'_i=Z_{i,j(i)}$, et on trouve un
  1096	plongement comme annoncé dans \hyperref[IX.prop2.14]{2.14}. Il reste à prouver le
  1097	
  1098	\begin{enonce*}{Lemme 2.14.1}\etiquette[2.14.1]{IX.lem2.14.1}
  1099	Soient $X$ un schéma géométriquement unibranche (par exemple normal),
  1100	$i:U\rightarrow X$ une immersion ouverte quasi-compacte. Alors
  1101	l'adhérence $Z$\pageoriginale de $U$, munie de la structure réduite
  1102	induite, est également géométriquement unibranche, plus précisément,
  1103	pour tout $z\in Z$, on a
  1104	$𝒪_{Z,z}\xleftarrow{\sim}𝒪_{X_{\red},z}$. Supposons d'autre part
  1105	$i$ dominant, et soit $C$ un ensemble, $C_U$ le faisceau constant sur
  1106	$U$ défini par $C$. Alors l'homomorphisme naturel $C_X\rightarrow
  1107	i_{\ast}(C_U)$ est un isomorphisme. Mêmes conclusions si on remplace
  1108	l'immersion ouverte $i$ par une inclusion $i:x\rightarrow X$, où $x$
  1109	est un point maximal de $X$.
  1110	\end{enonce*}
  1111	
  1112	Utilisant (\EGA IV 2.3.11), on est ramené pour la première assertion au
  1113	cas où $X$ est local, de point fermé $z$, mais alors $X_{\text{
  1114	réd}}$ est intègre et $U$ non vide, donc $Z=X_{\text{ réd}}$ et
  1115	l'assertion est évidente (il suffisait donc, au lieu de $X$
  1116	géométriquement unibranche, de supposer ici que les anneaux locaux de
  1117	$X_{\text{réd}}$ sont intègres). Pour la deuxième assertion,
  1118	regardant fibre par fibre, on se ramène grâce à (VIII 5.3) \loccit au
  1119	cas où $X$ est strictement local, et à prouver alors que
  1120	$H^0(X,C_X)\rightarrow H^0(U,C_U)$ est bijectif. Or ici encore, $X$
  1121	est irréductible, et $U$ en est un ouvert non vide donc irréductible,
  1122	et les deux membres de l'application précédente s'identifient à $C$,
  1123	d'où la conclusion voulue. (On a utilisé ici la notion « géométriquement unibranche » sous la forme de (\EGA IV 18.8.14),
  1124	comme signifiant que les anneaux localisés stricts de $X_{\text{
  1125	réd}}$ sont intègres). Le cas de l'inclusion d'un point maximal
  1126	(énoncé ici pour la commodité de références futures) se traite
  1127	exactement de la même façon.
  1128	
  1129	\begin{enonce*}[remark]{Remarques 2.14.2}\etiquette[2.14.2]{IX.rems2.14.2}
  1130	Lorsque $X$ est noethérien, la démonstration donnée (ou une déduction
  1131	immédiate) montre que dans \hyperref[IX.prop2.14]{2.14} (ii), on peut supposer les $X'_i$
  1132	intègres; si de plus on suppose $X$ universellement japonais, \ie les
  1133	anneaux de ses ouverts affines universellement japonais (\EGA ${\rm 0_{IV}}$
  1134	23.1.1, IV 7.7.2), on peut de plus supposer les $X'_i$ normaux.
  1135	\end{enonce*}
  1136	
  1137	\section{Théories de Kummer et d'Artin-Schreier}\etiquette[3]{IX.sec3}\pageoriginale
  1138	
  1139	\setcounter{subsection}{-1}
  1140	\subsection{}\etiquette[3.0]{IX.subsec3.0}
  1141	On va noter $(\GG_m)_X$ le faisceau représenté par le
  1142	groupe multiplicatif $\Spec𝒪_X[t,t^{-1}]$ sur $X$. Donc pour
  1143	$U\rightarrow X$ étale, les sections de $(\GG_m)_X$
  1144	sur $U$ sont les éléments inversibles $\Gamma(U,𝒪^{\ast}_U)$
  1145	de $\Gamma(U,𝒪_U)$. Soit $n\in \NN$. Le
  1146	noyau de la puissance $n$-ième dans $(\GG_m)_X$ est
  1147	le « faisceau des racines $n$-ièmes de l'unité », noté
  1148	$(𝛍_n)_X$ est le « faisceau des racines
  1149	$n$-ièmes de l'unité », noté $(𝛍_n)_X$. On a
  1150	donc une suite exacte
  1151	\begin{equation*}
  1152	0\rightarrow(𝛍_n)_X\rightarrow(\GG_m)_X\xrightarrow{n}(\GG_m)_X{\quoi}.\tag{3.1}
  1153	\end{equation*}\etiquette[3.1]{IX.eq3.1}
  1154	Si $n$ est inversible sur $X$ (c'est-à-dire, si $n$ est premier à
  1155	chaque caractéristique résiduelle de $X$), la puissance $n$-ième est
  1156	un morphisme surjectif \emph{pour les faisceaux}:
  1157	
  1158	\begin{enonce*}{Théorie de Kummer 3.2}\etiquette[3.2]{IX.Thr3.2}
  1159	On a la suite exacte
  1160	$$
  1161	0\rightarrow(𝛍_n)_X\rightarrow
  1162	(\GG_m)_X\xrightarrow{n}
  1163	(\GG_m)_X\rightarrow 0
  1164	$$
  1165	\end{enonce*}
  1166	\noindent
  1167	si $n\in \NN$ est inversible sur $X$.
  1168	
  1169	En effet, soit $u\in (\GG_m)_X(U)$, $U\rightarrow
  1170	X$ étale. Il faut trouver un recouvrement étale $\{U_i\rightarrow U\}$
  1171	de $U$ tel que $u$ induise une puissance $n$-ième sur chaque
  1172	$U_i$. Puisque $n$ est inversible sur $U$, l'équation
  1173	$$
  1174	T^n-u=0
  1175	$$
  1176	est « séparable »  sur $U$, \ie,
  1177	$$
  1178	U'=\Spec O_U[T]/(T^n-u)
  1179	$$
  1180	est étale au-dessus de $U$. Comme $U'\rightarrow U$ est surjectif, il
  1181	est couvrant et $u$ induit une puissance $n$-ième sur $U'$, on a fini.
  1182	
  1183	Signalons\pageoriginale que lorsque $n$ est inversible sur $X$,
  1184	$(𝛍_n)_X$ est un faisceau localement constant, en
  1185	fait localement isomorphe au faisceau
  1186	$(\ZZ/n\ZZ)_X$, et est constant
  1187	dans des cas importants. La théorie de Kummer donne donc des
  1188	renseignements sur la cohomologie de $X$ à valeurs dans certains
  1189	faisceaux constants. En fait, on a par définition
  1190	$$
  1191	H^0(X,(\GG_m)_X)=\Gamma(X,𝒪^{\ast}_X){\quoi},
  1192	$$
  1193	et pour la dimension 1 on a le
  1194	
  1195	\begin{enonce*}{Théorème 3.3}\etiquette[3.3]{IX.thm3.3}
  1196	(« {\rm théorème 90 de Hilbert} »): On a un isomorphisme
  1197	$$
  1198	H^1(X,(\GG_m)_X)\xrightarrow{\sim}\Pic X{\quoi},
  1199	$$
  1200	où $\Pic X$ est le groupe des classes de faisceaux inversibles sur $X$.
  1201	\end{enonce*}
  1202	
  1203	Démonstration. Cela veut dire que le morphisme canonique
  1204	$$
  1205	H^1(X_{\zar}, (\GG_m)_X)\longrightarrow
  1206	H^1(X_{\text{\rm et}},(\GG_m)_X)
  1207	$$
  1208	est bijectif, ce qui revient au même que de dire que
  1209	$$
  1210	R^1\epsilon_{\ast}(\GG_m)_X=O
  1211	$$
  1212	où $\epsilon_{\ast}$ est l'image directe pour le morphisme de sites
  1213	évident $\epsilon:X_{\et}\rightarrow X_{\zar}$ (VIII 4). On est donc
  1214	réduit au cas où $X$ est affine. Comme on peut calculer $H^1$ par
  1215	\v{C}ech (V 2.5) et comme il suffit pour un $X$ quasi-compact de
  1216	prendre des recouvrements quasi-compacts, c'est une conséquence
  1217	immédiate de la théorie de descente pour les faisceaux (\cf FGA 190,
  1218	ou \SGA VIII 1).
  1219	
  1220	\setcounter{subsection}{3}
  1221	\setcounter{subsubsection}{0}
  1222	\subsubsection{}\etiquette[3.3.1]{IX.subsubsec3.3.1}
  1223	Supposons $X$ noethérien et sans composante immergée et soient $x_j$
  1224	les points maximaux de $X$. Soit $R_j$ l'anneau artinien local de $X$
  1225	en $x_j$, et $i_j:\Spec R_j\rightarrow X$ le morphisme
  1226	d'inclusion. On\pageoriginale a une injection canonique
  1227	$$
  1228	(\GG_m)_X\longrightarrow \prod_j
  1229	i_{j^{\ast}}(\GG_m)_{\Spec R_j}
  1230	$$
  1231	de faisceaux sur $X_{\text{\rm et}}$ (\resp $X_{\zar}$). Le conoyau
  1232	$D$ est appelé le \emph{faisceau dés diviseurs de Cartier} sur
  1233	$X_{\text{\rm et}}$ (\resp $X_{\zar}$)\footnote[1]{Dans \EGA IV 21,
  1234	on dit simplement « diviseurs »  au lieu de « diviseurs de
  1235	Cartier ».}.
  1236	
  1237	\begin{enonce*}{Corollaire 3.4}\etiquette[3.4]{IX.cor3.4}
  1238	Soit $X$ noethérien et sans composante immergée, soit
  1239	$\epsilon:X_{\text{\rm ét}}\rightarrow X_{\zar}$ le morphisme de sites
  1240	canonique, et soit $D$ le faisceau des diviseurs de Cartier sur
  1241	$X_{\text{\rm et}}$. Alors $\epsilon_{\ast}D$ est le faisceau
  1242	$D_{\zar}$ des diviseurs de Cartier sur $X_{zar}$, et on a en
  1243	particulier:
  1244	$$
  1245	H^0(X_{\text{\rm et}},D)\approx H^0(X_{\zar}, D_{\zar}){\quoi}.
  1246	$$
  1247	\end{enonce*}
  1248	
  1249	Démonstration. Si l'on applique $\epsilon_{\ast}$ à la suite exacte
  1250	$$
  1251	0\rightarrow (\GG_m)_X\longrightarrow \prod_j
  1252	i_{j^{\ast}}(\GG_m)_{\Spec R_j}\longrightarrow
  1253	D\longrightarrow 0{\quoi},
  1254	$$
  1255	on trouve une suite exacte puisque
  1256	$R^1\epsilon_{\ast}(\GG_m)_X=0$ (\hyperref[IX.thm3.3]{3.3}), d'où la
  1257	première assertion.
  1258	
  1259	Dans le cas où $X$ admet une caractéristique $p\neq 0$, le résultat
  1260	suivant remplacera parfois \hyperref[IX.Thr3.2]{3.2} pour l'étude des coefficients de
  1261	$p$-torsion:
  1262	
  1263	\begin{enonce*}{Théorie d'Artin-Schreier 3.5}\etiquette[3.5]{IX.Thr3.5}
  1264	On a la suite exacte
  1265	$$
  1266	0\rightarrow (\ZZ/p
  1267	\ZZ)_X\rightarrow
  1268	𝒪_X\xrightarrow{\mathfrak{p}}𝒪_X\rightarrow 0{\quoi},
  1269	$$
  1270	si $X$ est de caractéristique $p>0$, où $\mathfrak{p}$ est le morphisme
  1271	de groupes additifs défini par $\mathfrak{p}(f)=f^p-f$.
  1272	\end{enonce*}
  1273	
  1274	Le noyau de $\mathfrak{p}$ est le faisceau des sections de $O_X$ qui
  1275	sont localement dans le corps premier, donc est in faisceau
  1276	constant. Le morphisme $\mathfrak{p}$\pageoriginale est surjectif pour
  1277	la topologie étale parce que l'équation $T^p-T-a$ $(a\in
  1278	\Gamma(U,𝒪_U))$ est toujours « séparable »  en caractéristique
  1279	$p$, \ie définit un revêtement étale de $U$. Pour appliquer cette
  1280	suite exacte on utilisera bien entendu (VII 4.2).
  1281	
  1282	\begin{enonce*}{Proposition 3.6}\etiquette[3.6]{IX.prop3.6}
  1283	\begin{enumerate}
  1284	\renewcommand{\theenumi}{\roman{enumi}}
  1285	\renewcommand{\labelenumi}{(\theenumi)}
  1286	\item Soit $X$ un schéma et $i:x\rightarrow X$ l'inclusion d'un
  1287	point. Alors
  1288	$$
  1289	H^1(X,i_{\ast}(\ZZ)_x)=0
  1290	$$
  1291	où $(\ZZ)_x$ est le « faisceau constant » {\rm
  1292	(\cf 2.0)} de valeur $\ZZ$.
  1293	
  1294	\item Si $X$ est irréductible et si pour chaque point géométrique
  1295	$\bar{x}$ de $X$ le localisé strict $X_{\bar{x}}$ de $X$ en $\bar{x}$
  1296	est irréductible (on devrait dire « $X$ est géométriquement
  1297	unibranche »  {\rm \cf \EGA IV 18.8.14}) alors on a
  1298	$$
  1299	H^{i}(X,\ZZ_x)=0{\quoi}.
  1300	$$
  1301	\end{enumerate}
  1302	\end{enonce*}
  1303	
  1304	Démonstration.
  1305	\begin{enumerate}
  1306	\renewcommand{\theenumi}{\roman{enumi}}
  1307	\renewcommand{\labelenumi}{(\theenumi)}
  1308	\item De la suite spectrale de Leray
  1309	$$
  1310	E^{pq}_2=H^p(X,R^{q}i_{\ast}(\ZZ)_x)\Longrightarrow
  1311	H^{p+q}(x,(\ZZ)_x)
  1312	$$
  1313	on déduit une injection
  1314	$H^1(X,i_{\ast}(\ZZ)_x)\rightarrow
  1315	H^1(x,(\ZZ)_x)$. Or puisque $x=\Spec k(x)$ est le
  1316	spectre d'un corps, il est bien connu que
  1317	$H^1(x,(\ZZ)_{x})=0$ (cela résulte du fait que
  1318	$H^q(x,F)(q>0)$ est de torsion en tout cas (\cf VII 2.2 et CG), et des
  1319	suites exactes $0\rightarrow
  1320	\ZZ\xrightarrow{n}\ZZ\rightarrow
  1321	\ZZ/n\rightarrow 0$), d'où le résultat.
  1322	
  1323	\item Soit $i:x\rightarrow X$ l'inclusion du point générique de
  1324	$X$. On voit immédiatement que sous les conditions de (ii) le
  1325	morphisme canonique $(\ZZ)_X\rightarrow
  1326	i_{\ast}(\ZZ)_x$ est bijectif (la réciproque étant
  1327	d'ailleurs vraie également!), et (ii) est ainsi réduit à (i).
  1328	\end{enumerate}
  1329	
  1330	\begin{enonce*}[remark]{Remarques 3.7}\etiquette[3.7]{IX.rems3.7}
  1331	Dans\pageoriginale \hyperref[IX.prop3.6]{3.6} (i) et (ii) on peut remplacer
  1332	$\ZZ$ par
  1333	n'importe quel groupe abélien sans torsion. D'autre part, (ii) devient
  1334	faux si on y abandonne l'hypothèse que $X$ est géométriquement
  1335	unibranche, comme on voit par exemple en prenant pour $X$ une courbe
  1336	algébrique sur un corps $K$, admettant un point double ordinaire
  1337	(faire le calcul de $H^1(X,\ZZ)$ dans ce cas!).
  1338	\end{enonce*}
  1339	
  1340	\section{Cas d'une courbe algébrique}\etiquette[4]{IX.sec4}
  1341	
  1342	\setcounter{subsection}{-1}
  1343	\subsection{}\etiquette[4.0]{IX.subsec4.0}
  1344	Soit $X$ un schéma noethérien de dimension 1. Soient $R(X)$ l'anneau
  1345	des fonctions rationnelles sur $X$, et $i:\Spec R(X)\rightarrow X$ le
  1346	morphisme d'inclusion. L'anneau $R(X)$ est artinien, produit des
  1347	anneaux locaux de $X$ aux points maximaux, et on peut utiliser les
  1348	résultats de (VIII 2).
  1349	
  1350	Soit $F$ un faisceau abélien sur $\Spec R(X)$ et considérons les
  1351	faisceaux $R^qi_{\ast}F$, $q>0$, sur $X$. Il résulte aussitôt de VIII
  1352	5.2 que la fibre de $R^qi_{\ast}F$ en un point géométrique au-dessus
  1353	d'un point maximal de $X$ est nulle. Puisque $X$ est de dimension 1,
  1354	un tel faisceau est de nature très spéciale:
  1355	
  1356	\begin{enonce*}{Lemme 4.1}\etiquette[4.1]{IX.lem4.1}
  1357	Soient $X$ noethérien et $F$ un faisceau sur $X$. Les conditions
  1358	suivantes sont équivalentes (on appellera un faisceau satisfaisant à
  1359	ces conditions un « skyscraper sheaf »):
  1360	\begin{enumerate}
  1361	\renewcommand{\theenumi}{\roman{enumi}}
  1362	\renewcommand{\labelenumi}{(\theenumi)}
  1363	\item Le fibre de $F$, en tout point géométrique $\bar{x}$ de $X$ qui
  1364	est au-dessus d'un point non-fermé $x\in X$, est nulle.
  1365	
  1366	\item Chaque section $s\in F(X')$ ($X'\rightarrow X$ étale de type
  1367	fini) est nulle, sauf en un nombre fini de points fermés de $X$.
  1368	
  1369	\item On a $F\approx ⨁_x i_{x^{\ast}}i_{\ast}{}^{\ast}F$, où $x$
  1370	parcourt l'ensemble des points fermés de $X$ et où $i_x:x\rightarrow
  1371	X$ est l'inclusion.
  1372	\end{enumerate}
  1373	\end{enonce*}
  1374	
  1375	Démonstration.\pageoriginale Supposons que (i) soit vrai, et soit
  1376	$z\in F(X')$ ($X'\rightarrow X$ étale de type fini). Alors le support
  1377	de $z$ (VIII 6.6) est une partie fermée de $X'$ formée de points
  1378	\emph{fermés}, donc finie ($X'$ étant noethérien), d'où (i)
  1379	$\Rightarrow$ (ii). L'implication (ii) $\Rightarrow$ (i) résulte du
  1380	calcul des fibres (VIII 3.9), et (iii) implique (i), car pour toute
  1381	famille $(G_x)$ de faisceaux sur les points fermés $x$ de $X$, le
  1382	faisceau $F=\oplus i_{x^{\ast}}(G_x)$ sur $X$ satisfait (i), comme il
  1383	résulte de fait que les foncteurs fibres commutent aux sommes
  1384	directes, et que le support de $i_{x^{\ast}}(G_x)$ est évidemment
  1385	contenu dans $\{x\}$.
  1386	
  1387	Il reste à démontrer que (ii) implique (iii). Considérons le morphisme
  1388	canonique
  1389	$$
  1390	F\longrightarrow \prod_x i_{x^{\ast}}i_{x}{}^\ast F{\quoi},
  1391	$$
  1392	$x$ parcourant l'ensemble des points fermés de $X$. D'après (ii), il
  1393	se factorise par
  1394	$$
  1395	F\longrightarrow \bigoplus_x i_{x^{\ast}}i_x{}^\ast F{\quoi},
  1396	$$
  1397	ce qui donne le morphisme cherché. Nous laissons au lecteur la
  1398	vérification du fait qu'il est bijectif.
  1399	
  1400	\setcounter{subsection}{1}
  1401	\subsubsection{}\etiquette[4.1.1]{IX.subsubsec4.1.1}
  1402	Appliquons le lemme au faisceau $R^qi_{\ast}F$ introduit plus haut: On
  1403	a
  1404	$$
  1405	R^qi_{\ast}F\xrightarrow{\sim} \bigoplus_x i_{x\ast} i_x \ast
  1406	(R^qi_{\ast}F)
  1407	$$
  1408	($x$ parcourant l'ensemble des points fermés de $X$). Or d'après (VIII
  1409	5) la fibre de $R^qi_{\ast}F$ en un point géométrique $\bar{x}$ de $X$
  1410	est
  1411	$$
  1412	(R^qi_{\ast}F)_{\bar{x}}\simeq H^q(\Spec R(X_{\bar{x}}), F)
  1413	$$
  1414	où $R(X_{\bar{x}})\simeq R(X)\otimes_{O_X}O_{X_{\bar{x}}}$,
  1415	$X_{\bar{x}}$ étant le localisé strict de $X$ en $\bar{x}$  (VIII
  1416	4),\pageoriginale et où on dénote aussi par $F$ le faisceau induit par
  1417	$F$ sur $\Spec R(X_{\bar{x}})$. Supposons maintenant que le corps
  1418	résiduel de tout point fermé $x$ de $X$ soit séparablement clos. On
  1419	aura donc (\cf aussi (VI. 1))
  1420	\begin{align*}
  1421	H^p(X,R^qi_{\ast}F) &\simeq H^p(X,\bigoplus_x i_{x^{\ast}}i_x\ast
  1422	(R^qi_{\ast}F))\\
  1423	&\simeq \bigoplus_x H^p(X, i_{x^{\ast}}i_x\ast(R^q i_{\ast}F))\\
  1424	&\simeq \bigoplus_x H^p(x,i_x\ast (R^q i_{\ast}F))\\
  1425	&\simeq \begin{cases}
  1426	         \displaystyle{\bigoplus_x} H^q(\Spec R(X_x),F) &\mbox{ si } p=0\\
  1427	          0                            &\mbox{ si } p>0{\quoi}.
  1428	         \end{cases}
  1429	\end{align*}
  1430	On trouve le
  1431	
  1432	\begin{enonce*}{Corollaire 4.2}\etiquette[4.2]{IX.cor4.2}
  1433	Soient $X$ un schéma noethérien de dimension $1$, tel que pour tout
  1434	point fermé $x$ de $X$, le corps résiduel $k(x)$ soit séparablement
  1435	clos. Soit $F$ un faisceau abélien, sur $\Spec R(X)$ (où $R(X)=$
  1436	anneau des fonctions rationnelles sur $X$). Considérons la suite
  1437	spectrale de Leray (\quad).
  1438	$$
  1439	E^{pq}_2=H^p(X,R^q i_{\ast}F)\Longrightarrow H^{p+q}(\Spec R(X),
  1440	F){\quoi}.
  1441	$$
  1442	On a
  1443	\begin{align*}
  1444	E^{pq} &= 0\mbox{ si } p \mbox{ et } q> 0{\quoi},\\
  1445	E^{0q}_2 &= \bigoplus_x H^q(\Spec R(X_x),F)\mbox{ si } q>0{\quoi},
  1446	\end{align*}
  1447	où $x$ parcourt l'ensemble des points fermés de $X$, et où $X_x$ est
  1448	le localisé strict de $X$ au point $x$.
  1449	
  1450	La\pageoriginale suite spectrale donne donc une suite exacte
  1451	\begin{align*}
  1452	& 0\longrightarrow H^1(X,i_{\ast}F)\longrightarrow H^1(\Spec R(X),
  1453	F)\longrightarrow \bigoplus_x H^1(\Spec R(X_x),F)\\
  1454	& \longrightarrow H^2(X,i_{\ast}F)\longrightarrow H^2(\Spec
  1455	R(X),F)\longrightarrow \bigoplus_x H^2(\Spec R(X_x),F)\\
  1456	&\longrightarrow\ldots
  1457	\end{align*}
  1458	\end{enonce*}
  1459	
  1460	\setcounter{subsection}{2}
  1461	\subsection{}\etiquette[4.3]{IX.subsec4.3}
  1462	Supposons maintenant que de plus la $\ell$-dimension cohomologique de
  1463	$\Spec R(X)$, pour un nombre premier $\ell$ donné, soit au plus égale
  1464	à 1 (c'est-à-dire qu'on ait $H^q(\Spec R(X), F)=0$, $q>1$, pour chaque
  1465	faisceau $F$ de $\ell$-torsion sur $\Spec R(X)$). Soient
  1466	$K^1,\ldots,K^m$ les corps résiduels de l'anneau artinien $R(X)$, \ie
  1467	les corps résiduels des points maximaux de $X$. Tenant compte du
  1468	dictionnaire (VIII 2) la $\ell$-dimension cohomologique de $\Spec
  1469	R(X)$ est
  1470	$$
  1471	\cd_{\ell}(\Spec R(X))=\mathop{\sup}_i \{\cd_{\ell}
  1472	(G(\overline{K}^{i}/K^{i}))\}{\quoi},
  1473	$$
  1474	où $\cd_{\ell}(G)$ est la $\ell$-dimension cohomologique du groupe
  1475	profini $G$ [\cf CG].
  1476	
  1477	Chaque corps résiduel $K_x$ de $R(\bar{X}(x))$, pour un localisé
  1478	strict $\bar{X}(x)$ de $X$, s'identifie évidemment à une extension
  1479	séparable (infinie) d'un des $K^i$, d'où
  1480	$$
  1481	\cd_{\ell}(G(\bar{K}_x/K_x))\leq \cd_{\ell} (G(\bar{K}^i/K^i))
  1482	$$
  1483	d'après (CG II 4.1, Prop. 10). Donc la $\ell$-dimension cohomologique
  1484	$\cd_{\ell}(\Spec R(\bar{X}(x)))$ est aussi $\leq 1$.
  1485	
  1486	En appliquant la suite exacte de \hyperref[IX.cor4.2]{4.2} à un $F$ de $\ell$-torsion, on
  1487	trouve
  1488	\begin{equation*}
  1489	\begin{cases}
  1490	 0\rightarrow H^1(X,i_{\ast}F)\rightarrow H^1(\Spec
  1491	R(X), F)\rightarrow\\
  1492	\qquad \displaystyle{\bigoplus_x} H^1(\Spec R(\bar{x}(x)),F)\rightarrow
  1493	H^2(X,i_{\ast}F)\rightarrow 0{\quoi},\\
  1494	 H^q(X,i_{\ast}F)=0\mbox{ pour } q> 2{\quoi}.
  1495	\end{cases}\tag{4.3.1}
  1496	\end{equation*}\etiquette[4.3.1]{IX.eq4.3.1}
  1497	
  1498	Remarquons\pageoriginale que l'hypothèse que $\cd_{\ell}(\Spec R(X))$
  1499	soit $\leq 1$
  1500	est satisfait si $X$ est une courbe algébrique sur un corps
  1501	séparablement clos. (\emph{Théorème de Tsen} [6, th.~\hyperref[IX.subsec4.3]{4.3}]). Des suites
  1502	exactes analogues ont été étudiées dans ce cas par $\Ogg[2]$ et
  1503	\v{S}afarevic.
  1504	
  1505	\subsection{}\etiquette[4.4]{IX.subsec4.4}
  1506	Supposons que $\ell$ soit inversible sur $X$ et prenons
  1507	$F=(\GG_m) R(X)$. Appliquant \hyperref[IX.thm3.3]{3.3} et [CG,
  1508	chap. I. prop. 12], on trouve, toujours sous l'hypothèse que
  1509	$\cd_{\ell}(\Spec R(X))\leq 1$.
  1510	\begin{equation*}
  1511	\begin{cases}
  1512	H^1(\Spec R(X), \GG_{m})=0\\
  1513	H^q (\Spec R(X), \GG_{m})\; \displaystyle{
  1514	\mathop{=}_{\ell} }\; 0\quad q\geq 2{\quoi},
  1515	\end{cases}\tag{4.4.1}
  1516	\end{equation*}\etiquette[4.4.1]{IX.eq4.4.1}
  1517	où le $\ell$ sous l'égalité veut dire que le groupe du premier membre,
  1518	qui est de torsion en tout cas, n'a pas de $\ell$-torsion.
  1519	
  1520	On a le même résultat si on remplace $R(X)$ par
  1521	$R(\bar{X}(x))$. Puisque le faisceau $\GG_m$ sur
  1522	$\Spec R(X)$ induit évidemment de faisceau $\GG_m$ sur
  1523	$\Spec R(\bar{X}(x))$, on déduit de \hyperref[IX.cor4.2]{4.2}
  1524	\begin{equation*}
  1525	\begin{cases}
  1526	H^1 (X, i_{\ast}(\GG_m)_{R(X)})=0\\
  1527	H^q (X, i_{\ast}(\GG_m)_{R(X)})\;
  1528	\displaystyle{\mathop{=}_{\ell}}\; 0,\mbox{ si } q> 1{\quoi}.
  1529	\end{cases}\tag{4.4.2}
  1530	\end{equation*}\etiquette[4.4.2]{IX.eq4.4.2}
  1531	Or pour tout faisceau $F$ sur $X$, on voit grâce à \hyperref[IX.lem4.1]{4.1} que le noyau
  1532	$P$ et conoyau $Q$ de l'homomorphisme canonique $F\rightarrow
  1533	i_{\ast}i^{\ast}F$ sont des « faisceaux gratte-ciel », donc (comme
  1534	les corps résiduels en les points fermés de $X$ sont supposés
  1535	séparablement clos) on a $H^{i}(X,P)=H^{i}(X,Q)=0$ pour $i>0$, ce qui
  1536	implique aussitôt, par la suite exacte de cohomologie, que
  1537	$$
  1538	H^{i}(X,F)\longrightarrow H^{i}(X,i_{\ast}i^{\ast}F)
  1539	$$
  1540	est\pageoriginale un isomorphisme pour $i\geq 2$. Compte tenu de \eqhyperref[IX.eq4.4.2]{4.4.2}
  1541	on en conclut la relation
  1542	\begin{equation*}
  1543	H^{i}(X,\GG_m)\displaystyle{\mathop{\;=\;}_{\ell}}0\mbox{
  1544	pour } i\geq 2{\quoi},\tag{4.5}
  1545	\end{equation*}\etiquette[4.5]{IX.eq4.5}
  1546	valable lorsque $X$ est un préschéma noethérien de dimension 1,
  1547	satisfaisant à $\cd_{\ell}(R(X))\leq 1$, et dont les corps résiduels
  1548	en les points fermés sont séparablement clos. Compte tenu de la
  1549	théorie de Kummer (\hyperref[IX.Thr3.2]{3.2}) et de (\hyperref[IX.thm3.3]{3.3}) on en conclut le
  1550	
  1551	\begin{enonce*}{Théorème 4.6}\etiquette[4.6]{IX.thm4.6}
  1552	Soient $X$ noethérien de dimension $1$, $n$ un entier qui est
  1553	inversible sur $X$, supposons que $\cd_{\ell}(R(X))\leq 1$ pour chaque
  1554	nombre premier $\ell$ qui divise $n$, et que pour tout point fermé $x$
  1555	de $X$, le corps résiduel $k(x)$ soit séparablement clos. Alors on a
  1556	$$
  1557	H^q(X, (𝛍_n)_X)=0\mbox{ si } q> 2{\quoi},
  1558	$$
  1559	et la cohomologie pour $q\leq 2$ est donnée par la suite exacte
  1560	\begin{align*}
  1561	0\rightarrow H^0(X,(𝛍_n)_X) &\rightarrow \Gamma
  1562	(X,𝒪^{\ast}_X)\xrightarrow{n} Γ(X,𝒪^{\ast}_X)\rightarrow\\
  1563	H^1(X,(𝛍_n)_X) &\rightarrow \Pic X\xrightarrow{n}
  1564	\Pic X\rightarrow H^2 (X,(𝛍_n)_X)\rightarrow 0{\quoi}.
  1565	\end{align*}
  1566	\end{enonce*}
  1567	
  1568	\begin{enonce*}{Corollaire 4.7}\etiquette[4.7]{IX.cor4.7}
  1569	Soit $X$ une courbe complète connexe sur un corps séparablement clos
  1570	$k$, de caractéristique première à $n$. Alors on a
  1571	\begin{align*}
  1572	H^0(X,(𝛍_n)_X) &\simeq
  1573	𝛍_n(k){\quoi},\\
  1574	H^1(X,(𝛍_n)_X) &\simeq {}_n(\Pic X)=\mbox{ensemble des éléments de $\Pic X$ dont l'ordre divise $n$,}\\
  1575	H^2(X,(𝛍_n)_X) &\simeq (\Pic X)/n \simeq (\ZZ/n\ZZ)^c{\quoi},
  1576	\end{align*}
  1577	où\pageoriginale $c$ est l'ensemble des composantes irréductibles de $X$.
  1578	\end{enonce*}
  1579	
  1580	Pour la dernière assertion, « rappelons » que le « schéma de
  1581	Picard » [TDTE V] d'un tel $X/\Spec k$ admet une décomposition
  1582	\begin{equation*}
  1583	0\longrightarrow A\longrightarrow \sheaf{\Pic}_{X/k}\longrightarrow
  1584	\ZZ^c\longrightarrow 0\tag{4.8}
  1585	\end{equation*}\etiquette[4.8]{IX.eq4.8}
  1586	où $A$ est un groupe algébrique connexe sur $\Spec k$. Or
  1587	l'application « multiplication par $n$ »  dans un tel $A$ est
  1588	étale, donc surjectif pour les points à valeurs dans $k$. Il s'ensuit
  1589	qu'on a
  1590	$$
  1591	{}_n(\Pic X)\approx {}_n(A(k))
  1592	$$
  1593	et
  1594	$$
  1595	(\Pic X)/n\approx (\ZZ/n
  1596	\ZZ)^c{\quoi}.
  1597	$$
  1598	
  1599	Si $X$ est lisse sur $\Spec k$, de sorte que la jacobienne $A$ est une
  1600	variété abélienne, le groupe ${}_n(\Pic X)$ est isomorphe à
  1601	$(\ZZ/n)^{2g}$, où $g$ est le genre de $X$. Bien
  1602	entendu, on a aussi $𝛍_n(k)\approx
  1603	\ZZ/n \ZZ$, donc les rangs des
  1604	groupes de cohomologie en tant que $\ZZ/n$-modules
  1605	sont $1$, $2g$, $1$, ce qui est le résultat auquel on devait s'attendre.
  1606	
  1607	\section{La méthode de la trace}\etiquette[5]{IX.sec5}
  1608	
  1609	\subsection{}\etiquette[5.1]{IX.subsec5.1}
  1610	Soit $X$ un schéma, $f:X'\rightarrow X$ un revêtement étale, et $F$ un
  1611	faisceau sur $X$. Alors les formules d'adjonction donnent un morphisme
  1612	de \emph{restriction}
  1613	$$
  1614	F\xrightarrow{\res}f_{\ast}f^{\ast}F{\quoi}.
  1615	$$
  1616	Comme $f$ est fini, on a $H^q(X,f_{\ast}f^{\ast}F)\simeq H^q(X',
  1617	f^{\ast}F)$ d'où un morphisme, appelé également restriction:
  1618	\begin{equation*}
  1619	H^q(X,F)\xrightarrow{\res}H^q(X',f^{\ast} F){\quoi},\tag{5.1.1}
  1620	\end{equation*}\etiquette[5.1.1]{IX.eq5.1.1}
  1621	qui\pageoriginale est d'ailleurs le morphisme évident. Mais dans le
  1622	cas $f$ étale on a aussi un autre morphisme, appelé \emph{morphisme
  1623	trace}
  1624	\begin{equation*}
  1625	f_{\ast}f^{\ast} F\xrightarrow{\tr} F{\quoi},\tag{5.1.2}
  1626	\end{equation*}\etiquette[5.1.2]{IX.eq5.1.2}
  1627	qui donne des morphismes sur la cohomologie
  1628	\begin{equation*}
  1629	H^q(X', f^{\ast} F)\xrightarrow{\tr}H^q
  1630	(X,F){\quoi}.\tag{5.1.3}
  1631	\end{equation*}\etiquette[5.1.3]{IX.eq5.1.3}
  1632	Le morphisme trace est caractérisé par les deux propriétés suivantes:
  1633	\begin{enumerate}
  1634	\renewcommand{\theenumi}{\roman{enumi}}
  1635	\renewcommand{\labelenumi}{(\theenumi)}
  1636	\item Il est de caractère local pour la topologie étale.
  1637	
  1638	\item Si $X'$ est somme disjointe de $d$ copies de $X$, de sorte que
  1639	$f_{\ast}f^{\ast} F\simeq F^d$, la trace est l'application somme
  1640	$F^d\rightarrow F$.
  1641	\end{enumerate}
  1642	
  1643	L'unicité est triviale à partir de (i), (ii) puisque tout revêtement
  1644	étale est localement constant. Ayant l'unicité, l'existence s'ensuit
  1645	par l'argument usuel de descente.
  1646	
  1647	De (ii) on déduit la formule suivante:
  1648	\begin{equation*}
  1649	\tr\circ \res:F\longrightarrow F\mbox{ est la multiplication par le
  1650	degré local de $f$.}\tag{5.1.4}
  1651	\end{equation*}\etiquette[5.1.4]{IX.eq5.1.4}
  1652	Si le degré est une constante $d$, on en déduit que le morphisme
  1653	$\tr\circ \res: H^q(X,F)\rightarrow H^q(X,F)$ est la multiplication
  1654	par $d$.
  1655	
  1656	\begin{enonce*}{Corollaire 5.2}\etiquette[5.2]{IX.cor5.2}
  1657	Soit $f:X'\rightarrow X$ un revêtement étale de degré constant $d$. Si
  1658	$F$ est un faisceau de $r$-torsion sur $X$ avec $(r,d)=1$, alors
  1659	$H^q(X',f^{\ast} F)=0$ implique que $H^q(X,F)=0$.
  1660	\end{enonce*}
  1661	
  1662	En effet, la multiplication par $d$ induit un isomorphisme de $F$,
  1663	donc un\pageoriginale isomorphisme sur la cohomologie. Si ce dernier
  1664	isomorphisme se factorise par zéro, on a $H^q(X,F)=0$.
  1665	
  1666	\setcounter{subsection}{2}
  1667	\subsubsection{}\etiquette[5.2.1]{IX.subsubsec5.2.1}
  1668	Soient maintenant $i:U\rightarrow X$ une immersion, $f:U'\rightarrow
  1669	U$ un revêtement étale de degré $d$, et supposons $f$ induit par un
  1670	morphisme fini $g:X'\rightarrow X$ par changement de base, \ie on a
  1671	alors le diagramme cartésien
  1672	$$
  1673	\xymatrix@C=2cm@R=1.25cm{
  1674	U' \ar[r]^{i'}\ar[d]_{f}& X'\ar[d]^{g}\\
  1675	U  \ar[r]^{i}& X
  1676	}
  1677	$$
  1678	Soit $F$ un faisceau sur $U$. On a
  1679	$$
  1680	\xymatrix@C=2.25cm@M=.25cm{
  1681	F \ar[r]^{\txt{res}} \ar@/_1.5pc/[rr]_{d}& f_{\ast}f^{\ast}F
  1682	\ar[r]^{\txt{tr}}& F,
  1683	}
  1684	$$
  1685	donc, comme $i_{\ast}f_{\ast}=g_{\ast}i'_{\ast}$,
  1686	$$
  1687	\xymatrix@C=2.25cm{
  1688	i_{\ast}F \ar@/_2pc/[rr]_{d} \ar[r]^{\txt{res}}& g_{\ast}i' {}_{\ast}f^{\ast}F \ar[r]^{\txt{tr}}& i_{\ast}F,
  1689	}\leqno{(5.3)}
  1690	$$
  1691	et comme $i_{!}f_{\ast}\simeq g_{\ast}i'_{!}$ (résulte p.ex. de (VIII
  1692	5.5)),
  1693	$$
  1694	\xymatrix@C=2cm{
  1695	i_{!}F \ar@/_2pc/[rr]_{d} \ar[r]^{\txt{res}}& g_{\ast}i' {}_{!}f^{\ast}F \ar[r]^{\txt{tr}}& i_{!}F,
  1696	}\leqno{(5.4)}
  1697	$$\etiquette[5.4]{IX.eq5.4}
  1698	d'où
  1699	$$
  1700	\xymatrix@C=2cm{
  1701	H^{q}(X,i_{\ast} F) \ar@/_2pc/[rr]_{d'} \ar[r]^{\txt{res}}& H^{q}(X', i' _{\ast}f^{\ast}F) \ar[r]^{\txt{tr}}& H^q(X,i_{\ast}F)
  1702	}\leqno{(5.3')}
  1703	$$\label{5.3'}
  1704	et\pageoriginale
  1705	$$
  1706	\xymatrix@C=2cm{
  1707	H^{q}(X,i_{!} F) \ar[r]^{\txt{res}} \ar@/_2.5pc/[rr]_{d}& H^{q}(X', i' _{!}f^{\ast}F)
  1708	\ar[r]^{\txt{tr}}& H^q(X,i_{!}F) ;
  1709	}\leqno{(5.4')}
  1710	$$\label{5.4'}
  1711	On déduit de \eqhyperref[IX.eq5.4]{5.4} les corollaires suivants:
  1712	
  1713	\begin{enonce*}{Proposition 5.5}\etiquette[5.5]{IX.prop5.5}
  1714	Soit $X$ un schéma quasi-compact et quasi-séparé, $\ell \in
  1715	\PP$, et soit $H$ un foncteur semi-exact à valeurs
  1716	dans {\rm (Ab)}, défini sur la catégorie des $\ell$-faisceaux sur $X$,
  1717	et compatible avec les limites inductives filtrantes. Les assertions
  1718	suivantes sont équivalentes:
  1719	\begin{enumerate}
  1720	\renewcommand{\theenumi}{\roman{enumi}}
  1721	\renewcommand{\labelenumi}{(\theenumi)}
  1722	\item $H=0$.
  1723	
  1724	\item $H(f_{\ast}i_{!}(\ZZ/\ell))=0$ pour chaque
  1725	$f:Y\rightarrow X$ fini, et chaque ouvert $i:U\rightarrow Y$ de $Y$,
  1726	tel que $U$ soit de présentation finie sur $X$. Si $X$ est noethérien,
  1727	on peut même se borner à prendre $Y$ intègre.
  1728	\end{enumerate}
  1729	\end{enonce*}
  1730	
  1731	Démonstration. D'après \hyperref[IX.cor2.7.2]{2.7.2} et \hyperref[IX.prop2.5]{2.5} il
  1732	suffit de vérifier $H(F)=0$
  1733	pour $F$ de la forme $F=i$, $G$, où $i:U\rightarrow X$ est l'inclusion
  1734	d'un sous-schéma de $X$ de présentation finie sur $X$, et où $G$ est
  1735	un $\ell$ faisceau localement constant « isotrivial »  sur
  1736	$U$. Soit $U''\rightarrow U$ un revêtement principal, de groupe $g$;
  1737	tel que $g$ devienne constant sur $U''$. Soit $A$ le $\ell$-groupe des
  1738	valeurs de $G$ sur $U''$. Le groupe $g$ opère sur $A$, et cette
  1739	opération détermine la structure de $G$. Soit $\sheaf{h}$ un
  1740	$\ell$-sous-groupe de Sylow de $g$, et soit $U'=U''/\sheaf{h}$ le
  1741	revêtement de $U$ correspondant à $\sheaf{h}$. Le degré $d$ de $U'$ sur
  1742	$U$ est premier à $\ell$. Prenons un prolongement de $f:U'\rightarrow
  1743	U$ en un morphisme fini $g:X'\rightarrow X$, qui existe toujours (\EGA
  1744	IV 8.12.6), et soit $i':U'\rightarrow X'$ l'inclusion. En appliquant
  1745	\eqhyperref[IX.eq5.4]{5.4} et le fait que $d$ est premier à $\ell$, on se réduit
  1746	à démontrer
  1747	que $H(g_{\ast}i'_{!}f^{\ast} G)=0$.
  1748	
  1749	Or\pageoriginale il est bien connu que le seul groupe abélien de
  1750	$\ell$-torsion qui est simple pour une opération d'un $\ell$ groupe
  1751	$\sheaf{h}$ est $\ZZ/\ell$ avec opération triviale, et
  1752	il s'ensuit qu'il existe une filtration de $A$ en tant que
  1753	$\sheaf{h}$-module, dont les quotients successifs sont isomorphes à
  1754	$\ZZ/\ell$ avec opération triviale. Cette
  1755	filtration donne une filtration correspondante de $f^{\ast} G$, et il
  1756	suffit donc, pour démontrer que $H(g_{\ast}i'_{!} f^{\ast} G)=0$ de
  1757	démontrer que $H(g_{\ast}i'_{!}\ZZ/r)=0$, d'où le
  1758	résultat.
  1759	
  1760	\begin{enonce*}{Proposition 5.6}\etiquette[5.6]{IX.prop5.6}
  1761	Soient $X$ un schéma noethérien, $\ell \in \PP$, et
  1762	$\varphi$ une fonction à valeurs dans $\NN$,
  1763	définie sur l'ensemble des schémas intègres finis sur $X$. Supposons
  1764	que $\varphi(Y_1)\leq \varphi(Y_2)$ chaque fois qu'il existe un
  1765	$X$-morphisme $Y_1\rightarrow Y_2$, et que l'inégalité est stricte si
  1766	$Y_1$ est un sous-schéma fermé de $Y_2$, distinct de $Y_2$. Pour tout
  1767	$Y$ fini sur $X$, posons $(Y)=\Sup (Y_i)$, où les $Y_i$ sont les
  1768	composantes irréductibles de $Y$, munis de la structure induit
  1769	réduite. Alors les assertions suivantes sont équivalentes:
  1770	\begin{enumerate}
  1771	\renewcommand{\theenumi}{\roman{enumi}}
  1772	\renewcommand{\labelenumi}{(\theenumi)}
  1773	\item Pour chaque $Y$ fini sur $X$ et chaque $\ell$-faisceau $F$ sur
  1774	$Y$ on a $H^q(Y,F)=0$ si $q>\varphi(\Supp F)$.
  1775	
  1776	\item Pour chaque $Y$ intègre fini sur $X$ on a
  1777	$H^q(Y,\ZZ/\ell)=0$ si $q>\varphi(Y)$.
  1778	\end{enumerate}
  1779	\end{enonce*}
  1780	
  1781	Démonstration. Supposons que (ii) soit vrai. On peut appliquer la
  1782	proposition précédente à chaque $H^q$, tenant compte que la
  1783	cohomologie commute aux limites inductives et aux images directes par
  1784	morphismes finis, et on se réduit ainsi pour la vérification de (i) au
  1785	cas $Y$ intègre fini sur $X$, et $F$ de la forme
  1786	$i_{!}(\ZZ/\ell)_U$, où $i:U\rightarrow X$ est un
  1787	ouvert non-vide.
  1788	
  1789	Or on a une suite exacte
  1790	$$
  1791	0\longrightarrow i_{!}(\ZZ/\ell)_U\longrightarrow
  1792	(\ZZ/\ell)_Y\longrightarrow
  1793	(\ZZ/\ell)_Z\longrightarrow 0
  1794	$$
  1795	où $Z=Y-U$ donc $Z\neq Y$. Faisant une récurrence sur la valeur de
  1796	$\varphi$, nous pouvons\pageoriginale supposer que
  1797	$H^q(\ZZ, \ZZ/\ell)=0$ pour
  1798	$q>\varphi (Y)-1$, d'où $H^q(Y,i_{!}(\ZZ/\ell))=0$
  1799	pour $q>\varphi(Y)$ grâce à la suite exacte de cohomologie.
  1800	
  1801	\begin{enonce*}{Corollaire 5.7}\etiquette[5.7]{IX.cor5.7}
  1802	Soit $X$ une courbe algébrique sur un corps $k$ séparablement clos et
  1803	de caractéristique différente de $\ell$, $\ell\in
  1804	\PP$. Alors la cohomologie de $X$ dans un faisceau
  1805	$F$ de $\ell$-torsion est nulle en dimension $>2$. Si $X$ est affine,
  1806	la cohomologie est nulle en dimension $>1$.
  1807	\end{enonce*}
  1808	
  1809	Posons, pour $Y$ fini intègre sur $X$, $\varphi(Y)=2\dim Y$ (\resp
  1810	$\varphi(Y)=\dim Y$ si $X$, donc $Y$, est affine). D'après 7.6 il
  1811	suffit de démontrer que $H^q(Y,\ZZ/\ell)=0$ si
  1812	$q>\varphi(Y)$, or c'est trivial si $\dim Y=0$. Il reste donc à
  1813	démontrer que $H^q(Y,\ZZ/\ell)=0$ si $q>2$ (\resp
  1814	si $q>1$ dans le cas affine) pour une courbe intègre $Y$. Puisque
  1815	$(\ZZ/\ell)_Y\simeq (𝛍_{\ell})_Y$, c'est
  1816	conséquence de \hyperref[IX.thm4.6]{4.6} pour $q>2$.
  1817	
  1818	Supposons $X$ affine. Puisque les extensions radicielles n'affectent
  1819	pas la cohomologie (VIII 1.1), on se réduit au cas où $k$ est
  1820	algébriquement clos. Soit $f:X'\rightarrow X$ la normalisation de $X$,
  1821	qui est un morphisme fini et un isomorphisme en dehors d'un ensemble
  1822	fini de points, disons $S$. On a une suite exacte de faisceaux sur $X$
  1823	$$
  1824	0\longrightarrow (\ZZ/n)_X\longrightarrow
  1825	f_{\ast}(\ZZ/n)_{X'}\longrightarrow \epsilon
  1826	\longrightarrow 0{\quoi},
  1827	$$
  1828	où $\epsilon$ est concentré sur $S$, donc où $H^q(X,\epsilon)=0$ pour
  1829	$q>0$. On se ramène par (VIII 5.6) à démontrer que
  1830	$$
  1831	H^q(X,f_{\ast}(\ZZ/n)_{X'})\simeq
  1832	H^q(X',\ZZ/n)=0
  1833	$$
  1834	si $q>1$, c'est-à-dire, au cas $X$ normale. On peut aussi évidemment
  1835	supposer $X$ connexe. Soit $i:X\rightarrow \bar{X}$ une immersion
  1836	ouverte avec $\bar{X}$ complète et non-singulière. Puisque $X$ est
  1837	affine, il manque au moins un point de $\bar{X}$, et on voit
  1838	immédiatement que par suite le morphisme
  1839	$$
  1840	A(k)\longrightarrow \Pic X
  1841	$$
  1842	est\pageoriginale surjectif, où $A$ est la Jacobienne de $\bar{X}$
  1843	(\cf \eqhyperref[IX.eq4.8]{4.8}). Puisque $A(k)$ est divisible par $\ell\neq \car k$, il en
  1844	est de même de $\Pic X$, d'où $H^2(X,𝛍_{\ell})=0$
  1845	par \hyperref[IX.thm4.6]{4.6}, \cqfd
  1846	
  1847	\begin{enonce*}[remark]{Remarque 5.8}\etiquette[5.8]{IX.rem5.8}
  1848	Signalons, pour référence ultérieure, que le raisonnement de \hyperref[IX.prop5.5]{5.5}
  1849	établit en fait le résultat suivant: Soient $X$ un schéma
  1850	quasi-compact et quasi-séparé, $\ell$ un nombre premier, $C$ une
  1851	sous-catégorie strictement pleine de la catégorie des faisceaux
  1852	abéliens de $\ell$-torsion constructibles, stable par facteurs directs
  1853	et par extensions. Supposons que pour chaque morphisme fini
  1854	$f:Y\rightarrow X$ et chaque ouvert $i:U\rightarrow Y$ de $Y$ qui est
  1855	de présentation finie sur $X$, on ait $f_{\ast}(i_{!}(𝐙/\ell𝐙)_U)\in C$. Alors $C$ contient tous les faisceaux
  1856	%% 𝒵 est ici ℤ [vérifier les macros]
  1857	constructibles de $\ell$-torsion.
  1858	\end{enonce*}
  1859	
  1860	\begin{thebibliography}{99}
  1861	\bibitem{1} Gabriel, P. Des Catégories Abéliennes,
  1862	Thèse, Paris (1962).
  1863	
  1864	\bibitem{2} Ogg, A. Cohomology of abelian varieties over
  1865	functions fields, Annals of Math. Vol. 76, No. 2 (1962).
  1866	
  1867	\bibitem{3} Deuring, M. Die Typen der Multiplikatorenringe elliptischer
  1868	Funktionenkörper, Hamb. Abh. Bol. 13 (1940) p. 197.
  1869	
  1870	\bibitem{4} Grothendieck, A. Fondements de la Géométrie
  1871	Algébrique, Extraits du Séminaire Bourbaki 1957-1962), Secrétariat
  1872	mathématique, 11, rue Pierre Curie, Paris 5ème. (\emph{Cité} FGA).
  1873	
  1874	\bibitem{5} Serre, J.P. Cohomologie Galoisienne, Lecture Notes
  1875	nᵒ 5, (Springer) (\emph{Cité} CG).
  1876	
  1877	\bibitem{6} Douady, A. Cohomologie des groupes compacts
  1878	totalement discontinus, Séminaire Bourbaki, exp. 189.
  1879	\end{thebibliography}
  1880	
  1881	\ifx\danslelivre\undefined
  1882	\end{document}
  1883	\fi
